% Proofs of the results of the main paper and of Section app:theory of this supplement (the full-cost floors, moved here from the main paper's Appendix A in revision v3; code/verify_theory.py reads Table tab:fullcost from sections/appendix_proofs.tex).

\section{Proofs}\label{app:proofs}

Throughout, $\ell=\lambda h(s)>0$, $\beta=1/\alpha>1$, $\kappa=\ell^{-\beta}$, $y=X-k$, $y_+=\max\{y,0\}$ and $C_\alpha=(1-\alpha)\alpha^{\alpha/(1-\alpha)}$. We use two elementary facts repeatedly.

\begin{enumerate}
\item[(F1)] For $c>0$ and $\beta>1$, the function $\Delta\mapsto c\Delta^{\beta}$ is strictly convex on $[0,\infty)$.
\item[(F2)] For $y>0$, $\ell>0$ and $\alpha\in(0,1)$, the function $\mathcal P(\sigma)=y\ell\sigma^{\alpha}-\sigma$ on $[0,\infty)$ is strictly concave and has the unique maximizer $\sigma^*=(\alpha\ell y)^{1/(1-\alpha)}$, with $\mathcal P(\sigma^*)=(1-\alpha)y\ell(\sigma^*)^{\alpha}=C_\alpha(\ell y)^{1/(1-\alpha)}$. Indeed, $\mathcal P''(\sigma)=\alpha(\alpha-1)y\ell\sigma^{\alpha-2}<0$ for $\sigma>0$, $\mathcal P'(\sigma)=\alpha y\ell\sigma^{\alpha-1}-1$ is positive near $0$ and tends to $-1$, and $\mathcal P'(\sigma^*)=0$ gives $\sigma^*$. At $\sigma^*$, $\sigma^*=\alpha y\ell(\sigma^*)^{\alpha}$, so $\mathcal P(\sigma^*)=(1-\alpha)y\ell(\sigma^*)^{\alpha}$. Finally, $\ell(\sigma^*)^{\alpha}=\ell^{1/(1-\alpha)}(\alpha y)^{\alpha/(1-\alpha)}$, which gives the last expression. For $y\le 0$, $\mathcal P(\sigma)\le-\sigma$, so the maximum is $0$, attained only at $\sigma=0$.
\end{enumerate}

\subsection*{Proof of \autoref{lem:timing}}

(a) In continuous time $K=kS$ gives $\dot K=k\dot S=kgS$, so investment per dollar of revenue is $kg$. In discrete time with $I_t=k(S_{t+1}-S_t)$, $I_t/S_t=k(S_{t+1}/S_t-1)=kg$ with $g$ the forward growth rate.

(b) With $I_t=k(S_t-S_{t-1})$, $I_t/S_t=k(1-S_{t-1}/S_t)=k\bigl(1-1/(1+g)\bigr)=kg/(1+g)$. Substituting the frontier spend $\sigma=\kappa(g-G)^{\beta}$ (\autoref{prop:frontier}(a), whose proof does not depend on the capital term) gives $\tilde m$. The difference is $kg-kg/(1+g)=kg^2/(1+g)\ge 0$ for $g>-1$; at $k=0.5$ and $g=0.3$ it is $0.045/1.3\approx 0.035$.

(c) Differentiating twice, $\tilde m''(g)=2k/(1+g)^3-\beta(\beta-1)\kappa(g-G)^{\beta-2}$ for $g>G$. If $\alpha<1/2$, then $\beta>2$ and the second term tends to $0$ as $g\downarrow G$, while the first tends to $2k/(1+G)^3>0$. By continuity, $\tilde m''>0$ on some interval $(G,G+\varepsilon)$, so $\tilde m$ is strictly convex there. The feasible set $\{(g,m):g\ge G,\ m\le\tilde m(g)\}$ is then not convex: for $G<g_1<g_2<G+\varepsilon$, the midpoint of $(g_1,\tilde m(g_1))$ and $(g_2,\tilde m(g_2))$ lies strictly above $\tilde m$ at the midpoint growth rate. For the example, let $\mathcal S(g)=g+0.5-1.25g/(1+g)-0.5(g-0.1)^{10/3}$. Then $\mathcal S'(G)=1-1.25/1.1^2<0$, so the corner $g=G$ is a local maximum with $\mathcal S(G)=0.1+0.5-0.125/1.1\approx 0.486$. Direct evaluation gives $\mathcal S'(0.5)\approx 0.248>0$, $\mathcal S'(0.7)\approx 0.061>0$ and $\mathcal S'(0.8)\approx -0.111<0$. So $\mathcal S$ first falls, then rises, then falls, and it has an interior local maximum in $(0.7,0.8)$, at $g\approx 0.739$ with $\mathcal S\approx 0.595$; \texttt{verify\_theory.py} (check 2d) locates it on a grid of step $10^{-5}$ and confirms that it is the only interior local maximum. Under \eqref{eq:fcf} with the same parameters, $X=1<k=1.25$, so the unique maximizer is $g=G$ by \autoref{prop:score}(c).

(d) If $\alpha=1/2$, then $\beta=2$ and $\tilde m''(g)=2k/(1+g)^3-2\kappa$, which is positive if and only if $(1+g)^3<k/\kappa$, that is, $g<(k/\kappa)^{1/3}-1$; the interval $(G,(k/\kappa)^{1/3}-1)$ is nonempty if and only if $k>\kappa(1+G)^3$. If $\alpha>1/2$, then $\beta<2$ and the second term of $\tilde m''$ tends to $\infty$ as $g\downarrow G$, so $\tilde m$ is strictly concave near $G$; for $g>G$, $\tilde m''(g)>0$ if and only if $2k(g-G)^{2-\beta}/(1+g)^3>\beta(\beta-1)\kappa$, and by continuity $\tilde m$ is then strictly convex on a neighborhood of that $g$, which makes the feasible set non-convex as in (c). The left side is positive for $g>G$ and does not depend on $\kappa$, so the condition holds for $\kappa$ small enough relative to $k$. In the asset-light example, $\kappa=0.86^{-1/0.3}\approx1.65$, and $\tilde m''(G+\Delta)=0$ at $\Delta\approx0.021$, with $\tilde m''>0$ below and $<0$ above (check 28b). \qed

\subsection*{Proof of \autoref{prop:frontier}}

(a) The map $\sigma\mapsto g=G+\ell\sigma^{\alpha}$ is a continuous strictly increasing bijection from $[0,\bar\sigma]\cap\mathbb R$ onto $I$, with inverse $\sigma(g)=((g-G)/\ell)^{1/\alpha}=\kappa(g-G)^{\beta}$. Every feasible $g$ is therefore produced by exactly one spend level, and the largest margin attainable at $g$ is $m_0-\sigma(g)-kg=m^*(g)$. Hence $\mathcal F=\{(g,m):g\in I,\ m\le m^*(g)\}$. This is the hypograph of a concave function (part (c)) over an interval, so it is convex. A point with $m<m^*(g)$ is dominated by $(g,m^*(g))$. Conversely, suppose $(g,m^*(g))$ were dominated by a feasible $(g',m')$, so that $g'\ge g$ and $m'\ge m^*(g)$ with one inequality strict. If $g'=g$, then $m'>m^*(g)$, which is infeasible. If $g'>g$, then $m'\le m^*(g')<m^*(g)$ because $m^*$ is strictly decreasing (part (c)), a contradiction.

(b) $m^*(G)=m_0-kG$. The right difference quotient is $[m^*(G+\Delta)-m^*(G)]/\Delta=-k-\kappa\Delta^{\beta-1}\to -k$ as $\Delta\downarrow 0$, because $\beta>1$.

(c) For $G\le g_1<g_2$ in $I$, $m^*(g_1)-m^*(g_2)=k(g_2-g_1)+\kappa[(g_2-G)^{\beta}-(g_1-G)^{\beta}]>0$, because $k\ge 0$, $\kappa>0$ and $\Delta\mapsto\Delta^{\beta}$ is strictly increasing on $[0,\infty)$. So $m^*$ is strictly decreasing on $I$, and for $g>G$, $m^{*\prime}(g)=-k-\beta\kappa(g-G)^{\beta-1}<-k\le 0$. Strict concavity on $I$ follows from (F1): $m^*$ is the sum of an affine function and $-\kappa(g-G)^{\beta}$, which is strictly concave on $[G,\infty)$. The second derivative $-\beta(\beta-1)\kappa(g-G)^{\beta-2}$ is negative for $g>G$. As $g\downarrow G$, $(g-G)^{\beta-2}$ tends to $0$ if $\beta>2$, equals $1$ if $\beta=2$, and tends to $\infty$ if $\beta<2$; with $\beta=1/\alpha$ this gives the three cases, and for $\beta=2$ the value is $-2\kappa$.

(d) With $\bar\sigma=\infty$, $I=[G,\infty)$. Since $\beta>1$, $\kappa(g-G)^{\beta}/g\to\infty$ as $g\to\infty$, so $m^*(g)/g=m_0/g-k-\kappa(g-G)^{\beta}/g\to-\infty$, which implies $m^*(g)\to-\infty$. Also $m^{*\prime}(g)=-k-\beta\kappa(g-G)^{\beta-1}\to-\infty$. \qed

\subsection*{Proof of \autoref{prop:score}}

(a) On the line $Xg+m=c$, $m=c-Xg$, which has slope $-X$. If $(g,m)\in\mathcal F$ with $m<m^*(g)$, then $(g,m^*(g))\in\mathcal F$ has a strictly higher score, so every maximizer lies on the frontier. On the frontier the score is $\mathcal R_X(g):=Xg+m^*(g)$, which is continuous and strictly concave on $I$. If $\bar\sigma<\infty$, $I$ is compact and a maximizer exists. If $\bar\sigma=\infty$, $\mathcal R_X(g)=g[X+m^*(g)/g]\to-\infty$ by \autoref{prop:frontier}(d), so a maximizer exists on a compact subinterval. Strict concavity implies that the maximizer is unique.

(b) Write $\mathcal R_X$ in terms of $\Delta=g-G$: $\mathcal R_X=m_0+yG+\mathcal S(\Delta)$ with $\mathcal S(\Delta)=y\Delta-\kappa\Delta^{\beta}$. For $y>0$, $\mathcal S'(\Delta)=y-\beta\kappa\Delta^{\beta-1}$ is positive near $0$ and strictly decreasing, so the unique stationary point solves $\beta\kappa(\Delta^*)^{\beta-1}=y$, that is,
\[
\Delta^*=\Bigl(\frac{y}{\beta\kappa}\Bigr)^{1/(\beta-1)}=\bigl(\alpha y\,\ell^{1/\alpha}\bigr)^{\alpha/(1-\alpha)}=\ell^{1/(1-\alpha)}(\alpha y)^{\alpha/(1-\alpha)},
\]
using $\beta-1=(1-\alpha)/\alpha$ and $\kappa^{-1}=\ell^{\beta}$. This is \eqref{eq:Dstar}. Since the cap does not bind, $\Delta^*\in(0,\bar g-G]$ and $g^*=G+\Delta^*$ is interior to $[G,\infty)$. The first-order condition $\mathcal S'(\Delta^*)=0$ reads $-m^{*\prime}(g^*)=k+\beta\kappa(\Delta^*)^{\beta-1}=k+y=X$, which is tangency. The spend producing $\Delta^*$ is $\sigma^*=\kappa(\Delta^*)^{\beta}=\Delta^*\cdot\kappa(\Delta^*)^{\beta-1}=\Delta^*\,y/\beta=\alpha y\Delta^*$. Substituting \eqref{eq:Dstar}, $\alpha y\Delta^*=\alpha y\,\ell^{1/(1-\alpha)}(\alpha y)^{\alpha/(1-\alpha)}=(\alpha\ell y)^{1/(1-\alpha)}$, which is \eqref{eq:sigstar} and coincides with the maximizer in (F2), as it must. Then $\mathcal S(\Delta^*)=y\Delta^*-\alpha y\Delta^*=(1-\alpha)y\Delta^*$, which gives the first form of \eqref{eq:Rstar}; the second follows from (F2), since $(1-\alpha)y\Delta^*=(1-\alpha)y\ell(\sigma^*)^{\alpha}$. Finally $m^*(g^*)=m_0-kg^*-\sigma^*=m_0-kg^*-\alpha y\Delta^*$.

(c) If $y\le 0$, then $\mathcal S(\Delta)=y\Delta-\kappa\Delta^{\beta}<0=\mathcal S(0)$ for every $\Delta>0$, so $g^*=G$ is the unique maximizer, and $R_X^*=m_0+yG$. The argument does not involve the cap.

(d) Let $G<g<\bar g$ and set $X(g):=-m^{*\prime}(g)=k+\beta\kappa(g-G)^{\beta-1}>k$. Then $\mathcal R_X$ with weight $X(g)$ is strictly concave with $\mathcal R_X'(g)=0$ at an interior point, so $g$ is its unique maximizer. For any other weight $X'\ne X(g)$, $\mathcal R_X'(g)=X'-X(g)\ne 0$ at the interior point $g$, so $g$ is not a maximizer. The point $g=G$ is a maximizer if and only if the right derivative $\mathcal R_X'(G^+)=X-k\le 0$, because $\mathcal R_X$ is concave. If $\bar\sigma=\infty$: for $X\le k$, $g^*(X)=G$; for $X>k$, $g^*(X)=G+\Delta^*(X)$ with $\Delta^*$ continuous, strictly increasing in $X$ and tending to $0$ as $X\downarrow k$. So $g^*$ is continuous and nondecreasing, and strictly increasing on $(k,\infty)$.

(e) $\Pi=\max_{\sigma\ge 0}\{y\ell\sigma^{\alpha}-\sigma\}$ is immediate from (F2) for $y>0$ and for $y\le 0$. Since $(X-k)G+(X-k)\ell\sigma^{\alpha}-\sigma+m_0=Xg+m_0-\sigma-kg$ along the frontier, $R_X^*=m_0+yG+\Pi$. \qed

\subsection*{Proof of Remarks~\ref{rem:cap} and \ref{rem:r40}}

For \autoref{rem:cap}, the score along the frontier as a function of spend is $m_0+yG+\mathcal P(\sigma)$ with $\mathcal P$ from (F2), restricted to $[0,\bar\sigma]$. Since $\mathcal P$ is strictly concave and increasing on $[0,\sigma^*]$, if $\bar\sigma<\sigma^*$ the maximum over $[0,\bar\sigma]$ is at $\bar\sigma$. When the cap binds strictly, $\bar\sigma<\sigma^*$ holds on a neighborhood of the parameters and the value is $m_0+yG+y\ell\bar\sigma^{\alpha}-\bar\sigma$, whose derivatives are those of \eqref{eq:envelope} with $\Delta^*$ replaced by $\ell\bar\sigma^{\alpha}$ and $\sigma^*$ by $\bar\sigma$. For \autoref{rem:r40}, $g\mapsto g+m^*(g)$ is the sum of a linear and a strictly concave function, so its upper level set $\{g:g+m^*(g)\ge 0.4\}$ is an interval. If the interval is nonempty, the maximum of $g+m^*(g)$ is at least $0.4$, so the maximizer belongs to it. \qed

\subsection*{Proof of \autoref{rem:general-phi}}

(a) The spend that produces growth $g\ge G$ is $\phi^{-1}((g-G)/\ell)$, which is convex in $g$ because $\phi$ is concave and strictly increasing, and strictly increasing. So $m^*(g)=m_0-kg-\phi^{-1}((g-G)/\ell)$ is concave and strictly decreasing, its hypograph over the feasible growth interval is convex, and the arguments of Lemmas~\ref{prop:frontier}(a) and \ref{prop:score}(a) apply unchanged: dominated points lie below the frontier, and a linear score is maximized on it. At an interior maximizer of the concave function $Xg+m^*(g)$ the derivative is zero, $m^{*\prime}(g^*)=-X$; strict concavity of $\phi$ makes $\phi^{-1}$ strictly convex, so the maximizer is unique and the argument of \autoref{prop:score}(d) applies. The proof of \autoref{prop:budget} uses only $m\le m_0-kg$, $g\ge G$ (true because $\phi\ge0$) and $\dot s/s=g-G$. \autoref{prop:industry}(iv) exhibits observationally equivalent isoelastic technologies, which belong to the general class. \autoref{prop:entry} has no growth technology. The envelope theorem gives $\partial R_X^*/\partial m_0=1$, $\partial R_X^*/\partial G=X-k$, $\partial R_X^*/\partial k=-g^*$ and $\partial R_X^*/\partial X=g^*$ wherever the maximizer is unique, because $R_X=m_0+(X-k)G+(X-k)\ell\phi(\sigma)-\sigma$ along the frontier.

(b) Along the frontier the score is $m_0+(X-k)G+(X-k)\ell\phi(\sigma)-\sigma$, so an interior optimum satisfies $(X-k)\ell\phi'(\sigma^*)=1$. Then $\sigma^*/\Delta^*=\sigma^*/(\ell\phi(\sigma^*))=(X-k)\sigma^*\phi'(\sigma^*)/\phi(\sigma^*)=(X-k)\varepsilon(\sigma^*)$. Concavity and $\phi(0)=0$ give $\sigma\phi'(\sigma)\le\phi(\sigma)$, so $\varepsilon\le1$, and strict monotonicity gives $\varepsilon>0$.

(c) The right derivative of $\phi^{-1}(\Delta/\ell)$ at $\Delta=0$ is $1/(\ell\phi'(0^+))$, which is zero if and only if $\phi'(0^+)=\infty$. The marginal score of the first dollar of spend is $(X-k)\ell\phi'(0^+)-1$; if it is not positive, concavity makes the score nonincreasing in spend, and the firm spends nothing.

(d) Each listed result uses the closed form \eqref{eq:Dstar}--\eqref{eq:Rstar}, which requires $\phi(\sigma)=\sigma^\alpha$. \qed

\subsection*{Proof of \autoref{prop:saturation}}

(a) By \autoref{prop:score}(b), (c) and (e), $R_X^*(s)=R_\infty+C_\alpha(\lambda h(s)y_+)^{1/(1-\alpha)}$ for every $X\ge 0$. If $y\le 0$, this is constant. If $y>0$, the function $\ell\mapsto C_\alpha(\ell y)^{1/(1-\alpha)}$ is strictly increasing, and $h(s)=1-s/\bar s$ is strictly decreasing, so $R_X^*$ is strictly decreasing in $s$. By the chain rule and (F2),
\[
\frac{dR_X^*}{ds}=\frac{C_\alpha}{1-\alpha}\,(\lambda y)^{1/(1-\alpha)}h^{\alpha/(1-\alpha)}\cdot\Bigl(-\frac{1}{\bar s}\Bigr)=-\frac{\Pi}{(1-\alpha)\,h\,\bar s}=-\frac{y\Delta^*}{\bar s-s},
\]
using $\Pi=(1-\alpha)y\Delta^*$ and $h\bar s=\bar s-s$.

(b) As $s\to\bar s$, $h(s)\to 0$, and $1/(1-\alpha)>0$, so the growth surplus tends to $0$ and $R_X^*(s)\to R_\infty$. For $X=1$ and $k=0$, $R_\infty=m_0+G$. Since $c_h\ge 0$ and $a\ge a_{\min}$, $m_0=1-c_h-c_c-f-a\le 1-c_c-f-a_{\min}\le 1$, which gives \eqref{eq:rule100}.

(c) By (a), $R_X^*(s)$ is nonincreasing in $s$ and continuous in $h$, so its supremum over $(0,\bar s)$ is its limit as $s\to 0$ ($h\to 1$), namely $R^0$. If $y>0$, it is strictly decreasing, so the supremum is not attained on the open interval; if $y\le0$, it is constant by (a), so the supremum is attained at every share.

(d) For $y>0$, $R_X^*(s)\ge\bar R$ if and only if $C_\alpha(\ell y)^{1/(1-\alpha)}\ge\bar R-R_\infty>0$. Because $x\mapsto x^{1-\alpha}$ is strictly increasing on $(0,\infty)$, this is equivalent to $\ell y\ge[(\bar R-R_\infty)/C_\alpha]^{1-\alpha}$, that is, $\ell\ge\Lambda(\bar R)$. Since $\ell=\lambda h(s)$ takes every value in $(0,\lambda)$ as $s$ ranges over $(0,\bar s)$, some share attains $\bar R$ if and only if $\lambda>\Lambda(\bar R)$. At $\ell=\Lambda(\bar R)$, $\Pi=\bar R-R_\infty$, so $\Delta^*=\Pi/[(1-\alpha)y]$ and $\sigma^*=\alpha y\Delta^*=\alpha\Pi/(1-\alpha)$, and the margin is $m^*(g^*)=m_0-kg^*-\sigma^*$ by \autoref{prop:score}(b). \qed

\subsection*{Proof of \autoref{prop:myopic}}

Write $\hat s=s/\bar s$, $h=1-\hat s$ and $\hat s_0=s_0/\bar s$, so that $h_0=1-\hat s_0$, and let $q:=1/(1-\alpha)$. Under the myopic policy, $\dot s/s=\Delta^*(\ell(s),X)=\Delta_0^*\,h^{q}$, because $\Delta^*=(\lambda h)^{q}(\alpha y)^{\alpha q}$. Hence $\dot h=-\dot{\hat s}=-\hat s\,\Delta_0^*h^{q}=-\Delta_0^*(1-h)h^{q}$. The right side is locally Lipschitz in $h$ on $(0,1)$, the unique solution stays in $(0,1)$ (the constant solutions $h\equiv 0$ and $h\equiv 1$ cannot be crossed), and $h$ is strictly decreasing. Thus $\hat s(t)$ is strictly increasing and bounded by $1$, and $\hat s(t)\ge\hat s_0$. Since $\dot h\le-\Delta_0^*\hat s_0h^{q}$, the function $\mathcal J:=h^{1-q}$ satisfies $\dot{\mathcal J}=(1-q)h^{-q}\dot h\ge(q-1)\Delta_0^*\hat s_0$, where $q-1=\alpha/(1-\alpha)>0$. Integrating, $h(t)^{-\alpha/(1-\alpha)}\ge h_0^{-\alpha/(1-\alpha)}+\tfrac{\alpha}{1-\alpha}\Delta_0^*\hat s_0t$, so $h(t)\to 0$, $s(t)\uparrow\bar s$, and
\[
R_X^*(s(t))-R_\infty=C_\alpha(\lambda y)^{q}h(t)^{q}\le C_\alpha(\lambda y)^{q}\Bigl[h_0^{-\alpha/(1-\alpha)}+\tfrac{\alpha}{1-\alpha}\Delta_0^*\hat s_0t\Bigr]^{-q(1-\alpha)/\alpha},
\]
with exponent $-q(1-\alpha)/\alpha=-1/\alpha$. This is \eqref{eq:myopic}. The left side is positive because $h(t)>0$, and it is decreasing in $t$ by \autoref{prop:saturation}(a). \qed

\subsection*{Proof of \autoref{prop:budget}}

(a) Let $\hat s=s/\bar s$ and $\Sigma(t):=\lambda\int_0^t\sigma(\tau)^{\alpha}d\tau<\infty$. The share equation $\dot s=s\,\lambda(1-s/\bar s)\sigma(t)^{\alpha}$ has the unique absolutely continuous solution $\hat s(t)=\bigl[1+e^{-\Sigma(t)}(1-\hat s_0)/\hat s_0\bigr]^{-1}$ (the logistic transform $\ln[\hat s/(1-\hat s)]$ has derivative $\lambda\sigma^{\alpha}$), so $\hat s(t)\in(0,1)$ and $s(t)<\bar s$ for all $t$. By \autoref{ass:market}, $\dot s/s=g-G=\Delta$ almost everywhere, and $\ln s$ is absolutely continuous, so $\int_0^T\Delta\,dt=\ln(s(T)/s_0)$. This is the only place where the functional form \eqref{eq:growth} is used; the rest of the argument needs only $\dot s/s=\Delta$ and $s\le\bar s$. Along any feasible path, $R_X=Xg+m\le Xg+m_0-\sigma-kg=R_\infty+y\Delta-\sigma\le R_\infty+y\Delta$, so $(R_X-R_\infty)^+\le y\Delta$ pointwise, because $y>0$ and $\Delta\ge0$: growth is at least $G$ whatever the spend. Integrating gives the first inequality in \eqref{eq:budget}; the second follows from $s(T)<\bar s$.

(b) By definition, $1+g_t=(1+G)s_{t+1}/s_t$. Along any feasible path, $R_{X,t}-R_\infty\le X g_t+m_0-kg_t-R_\infty=y(g_t-G)$, so $(R_{X,t}-R_\infty)^+\le y(g_t-G)$ because $g_t\ge G$. Therefore
\[
1+\frac{(R_{X,t}-R_\infty)^+}{y(1+G)}\le 1+\frac{g_t-G}{1+G}=\frac{1+g_t}{1+G}=\frac{s_{t+1}}{s_t}.
\]
Taking logarithms and summing over $t=0,\dots,T-1$ gives \eqref{eq:budget-discrete}, and $s_T\le\bar s$ gives the second inequality; $G>-1$ makes the logarithms well defined.

(c) With a moving intercept, $m(t)\le m_0(t)-\sigma(t)-kg(t)$, so $R_X(t)\le R_\infty(t)+y\Delta(t)-\sigma(t)$ at every date, and the arguments of (a) and (b) go through date by date with $R_\infty(t)$. \qed

\subsection*{Proof of \autoref{prop:budget-sharp}}

Fix $\underline\Delta>0$ and a target share $s_1\in(s_0,\bar s)$, and consider the policy that holds excess growth at $\underline\Delta$ until $s$ reaches $s_1$, at date $T_1=\ln(s_1/s_0)/\underline\Delta$, using spend $\sigma(t)=(\underline\Delta/(\lambda h(s(t))))^{1/\alpha}$ and frontier margins. Along the path $h(s(t))\ge h(s_1)>0$, so $\sigma(t)\le\bar\sigma_1:=(\underline\Delta/(\lambda h(s_1)))^{1/\alpha}$, which is finite and, for small $\underline\Delta$, below any positive financing cap. The excess score is $R_X-R_\infty=y\underline\Delta-\sigma(t)\ge y\underline\Delta-\bar\sigma_1$, which is positive for small $\underline\Delta$ because $\bar\sigma_1/\underline\Delta=\underline\Delta^{\beta-1}(\lambda h(s_1))^{-\beta}\to 0$. Hence
\[
\int_0^{T_1}(R_X-R_\infty)^+dt\ge T_1(y\underline\Delta-\bar\sigma_1)=y\ln\frac{s_1}{s_0}\Bigl(1-\frac{\bar\sigma_1}{y\underline\Delta}\Bigr)\longrightarrow y\ln\frac{s_1}{s_0}\quad(\underline\Delta\to 0).
\]
Letting $s_1\uparrow\bar s$ shows that the supremum is at least $y\ln(\bar s/s_0)$; by \autoref{prop:budget}(a) it is at most this value. \qed

\subsection*{Proof of \autoref{cor:transient}}

If $\mathcal T$ periods have $R_{X,t}\ge\bar R$, each contributes at least $\ln(1+(\bar R-R_\infty)/[y(1+G)])>0$ to the sum in \eqref{eq:budget-discrete}, and the other terms are nonnegative, so $\mathcal T\ln(1+(\bar R-R_\infty)/[y(1+G)])\le\ln(\bar s/s_0)$, which is \eqref{eq:transient}. In the example, $y=1$, $R_\infty=0.8$ and $1+G=1.2$. For $\bar R=10$, $\ln(1+9.2/1.2)=\ln 8.667\approx 2.159$; with $s_0=0.01$, $\ln(\bar s/s_0)=\ln 50\approx 3.912$ and the bound is $1.81$, so $\mathcal T(\bar R)\le 1$; with $s_0=10^{-5}$, $\ln(5\times 10^4)\approx 10.820$ and the bound is $5.01$, so $\mathcal T(\bar R)\le 5$. For $\bar R=1.4$ and $s_0=0.01$, $\ln(1+0.6/1.2)=\ln 1.5\approx 0.405$ and the bound is $9.65$, so $\mathcal T(\bar R)\le 9$.

(b) Let $\mathcal T_c$ be the measure of the set of dates with $R_X(t)\ge\bar R$. On that set $(R_X-R_\infty)^+\ge\bar R-R_\infty>0$, so \autoref{prop:budget}(a) gives $\mathcal T_c(\bar R-R_\infty)\le(X-k)\ln(\bar s/s_0)$. A firm at the static optimum for the weight $X$ at every date has $R_X(t)-R_\infty=\Pi(s(t))=(1-\alpha)(X-k)\Delta^*(s(t))$ by \autoref{prop:score}(e), so $\int(R_X-R_\infty)\,dt=(1-\alpha)(X-k)\int\Delta\,dt\le(1-\alpha)(X-k)\ln(\bar s/s_0)$, and the same argument gives the sharper bound. The asymmetry comes from the spend that excess growth requires: along any path $R_X-R_\infty\le(X-k)\Delta-\sigma$ with $\sigma\ge(\Delta/\lambda)^{1/\alpha}$, so fast spending converts headroom into score at a rate below $X-k$. In the example, $(X-k)\ln(\bar s/s_0)/(\bar R-R_\infty)=3.912/9.2\approx0.43$, and $0.7$ times that is $0.30$. \qed

\subsection*{Proof of \autoref{prop:envelope}}

By \autoref{prop:score}, $R_X^*=m_0+yG+C_\alpha(\lambda h\,y_+)^{1/(1-\alpha)}$. The function $y\mapsto y_+^{1/(1-\alpha)}$ is continuously differentiable on $\mathbb R$, with derivative $\frac{1}{1-\alpha}y_+^{\alpha/(1-\alpha)}$, which vanishes at $y=0$, because $1/(1-\alpha)>1$. Hence $R_X^*$ is continuously differentiable in all arguments. Write $\Pi(\ell,y):=C_\alpha(\ell y_+)^{1/(1-\alpha)}$ for the growth surplus. For $y>0$,
\[
\frac{\partial\Pi}{\partial y}=\frac{C_\alpha}{1-\alpha}\,\ell^{1/(1-\alpha)}y^{\alpha/(1-\alpha)}=\alpha^{\alpha/(1-\alpha)}\ell^{1/(1-\alpha)}y^{\alpha/(1-\alpha)}=\Delta^*,
\]
\[
\frac{\partial\Pi}{\partial\ell}=\frac{C_\alpha}{1-\alpha}\,y^{1/(1-\alpha)}\ell^{\alpha/(1-\alpha)}=\frac{y\Delta^*}{\ell},
\]
and both are zero for $y\le 0$, where $\Delta^*=0$. Since $y=X-k$ and $\ell=\lambda h$: $\partial R_X^*/\partial m_0=1$; $\partial R_X^*/\partial G=y$; $\partial R_X^*/\partial k=-G-\partial\Pi/\partial y=-(G+\Delta^*)=-g^*$; $\partial R_X^*/\partial X=G+\Delta^*=g^*$; and $\partial R_X^*/\partial\lambda=h\,\partial\Pi/\partial\ell=y\Delta^*/\lambda$. Using \eqref{eq:Dstar}, $\Delta^*=\ell(\sigma^*)^{\alpha}$, so $y\Delta^*/\lambda=y\,h\,(\sigma^*)^{\alpha}$, which is positive if and only if $y>0$. The derivatives with respect to $\gamma$, $a$ and $c_c$ follow from $m_0=1-c_h-c_c-f-a$ and $\gamma=1-c_h-c_c-f$.

Convexity in $k$: $\partial R_X^*/\partial k=-G-\Delta^*(y)$ with $y=X-k$, and $\Delta^*$ is nondecreasing in $y$ (strictly for $y>0$), hence nonincreasing in $k$; so $\partial R_X^*/\partial k$ is nondecreasing in $k$. Convexity in $\lambda$: $\Pi$ is proportional to $\lambda^{1/(1-\alpha)}$ with exponent above one. For the cross-partial, $\Delta^*=\ell^{1/(1-\alpha)}(\alpha y_+)^{\alpha/(1-\alpha)}$ is differentiable in $\lambda$ for every $y$, with $\partial\Delta^*/\partial\lambda=\Delta^*/[(1-\alpha)\lambda]$ (zero when $y\le 0$). Hence $\partial^2R_X^*/\partial\lambda\,\partial k=-\Delta^*/[(1-\alpha)\lambda]\le 0$; the mixed partial is continuous, so the order of differentiation does not matter. \qed

\subsection*{Proof of \autoref{cor:capital}}

Let $\mathcal R(k):=R_X^*$ at $a$ fixed. The total change is $\Delta R_X^*=da+\mathcal R(k+dk)-\mathcal R(k)$, because $R_X^*$ is linear in $a$ with slope $-1$. By \autoref{prop:envelope}, $\mathcal R$ is convex and differentiable with $\mathcal R'(k)=-g^*(k)$. Convexity gives $\mathcal R'(k)\,dk\le\mathcal R(k+dk)-\mathcal R(k)\le\mathcal R'(k+dk)\,dk$, which is \eqref{eq:kl-bracket}. If $da>g^*(k)dk$, the lower bound is positive. If $da\le g^*(k+dk)dk$, the upper bound is nonpositive. To first order, $\Delta R_X^*=da-g^*dk+o(dk)$. \qed

\subsection*{Proof of \autoref{cor:priorities}}

(i) $R_X^*=m_0+yG+\Pi$ and $R_\infty=m_0+yG$ are both linear in $a$ with slope $-1$, so removing $da$ raises each by $da$ exactly.

(ii) $R_X^*$ depends on $\lambda$ only through $\Pi=C_\alpha(\lambda hy)^{1/(1-\alpha)}$, so replacing $\lambda$ by $\lambda(1+x_\lambda)$ multiplies $\Pi$ by $(1+x_\lambda)^{1/(1-\alpha)}$ and changes $R_X^*$ by exactly $\Pi[(1+x_\lambda)^{1/(1-\alpha)}-1]$. To first order this is $\Pi x_\lambda/(1-\alpha)=y\Delta^*x_\lambda$, using $\Pi=(1-\alpha)y\Delta^*$, and $y\Delta^*=\sigma^*/\alpha$ by \eqref{eq:sigstar}. An improvement $x$ in spend efficiency replaces $\sigma$ by $\sigma/(1+x)$ in the growth technology, so $\lambda\sigma^\alpha$ becomes $\lambda(1+x)^\alpha\sigma^\alpha$: demand pull is multiplied by $(1+x)^\alpha$, and $\Pi$ by $(1+x)^{\alpha/(1-\alpha)}$. To first order the gain is $\Pi\,x\,\alpha/(1-\alpha)=\alpha y\Delta^*x=x\sigma^*$. Cutting a non-discretionary cost line of size $c_L$ by the proportion $x$ raises $R_X^*$ by $xc_L$ by (i), which equals $x\sigma^*$ when $c_L=\sigma^*$. $R_\infty$ does not depend on $\lambda$, and both gains tend to zero as $s\to\bar s$ because $\Pi\to 0$ (\autoref{prop:saturation}(b)).

(iii) With $g=G_i+\ell\sigma^{\alpha}$, the proof of \autoref{prop:score} applies with $G$ replaced by $G_i$ (\autoref{rem:assumptions}(b)), so $R_X^*=m_0+yG_i+\Pi$ and $R_\infty=m_0+yG_i$, both with derivative $y$ in $G_i$, which does not depend on $s$ or $g^*$.

(iv) By \autoref{lem:prices}(a), $d\gamma=-(1-\epsilon_p)du/p$ to first order, which is positive for $du<0$ and $\epsilon_p<1$, and $\partial R_X^*/\partial\gamma=1$. Under free entry without a scarce input, \autoref{prop:entry}(b) gives $L_{n^*}\to 0$ as $F\to 0$ for every $u$, and the limiting score $XG+(r-G)(k_u+a_c)/\hat c$ with $\hat c=u+r(k_u+a_c)$ depends on $u$ only through $(k_u+a_c)/\hat c$, and not at all if $k_u+a_c=0$.

(v) This is \autoref{cor:capital}; the first-order hurdle $g^*dk$ is increasing in $g^*$, and $g^*\ge G$. \qed

\subsection*{Proof of \autoref{prop:dynamic}}

{\sloppy (a) For any policy, the share path does not depend on $M_0$, and $S_t=s_tM_0e^{Gt}$. Hence $\int_0^Te^{-rt}m_tS_t\,dt=M_0\int_0^Te^{-(r-G)t}m_ts_t\,dt$, and taking the limit superior and the supremum over policies gives $V(S_0,M_0)=M_0\,v(s_0)$ with $v(s):=V(s,1)$. The policy $\sigma\equiv 0$ keeps $s_t=s_0$ and $m_t=m_0-kG$, which gives the lower bound. For any policy, $m_t\le m_0-kG$ because $\sigma_t\ge 0$, $g_t\ge G$ and $k\ge 0$, and $0<s_t<\bar s$ by the proof of \autoref{prop:budget}(a); so every integral is at most $(m_0-kG)\bar s/(r-G)$. Finally $V(S,M)=Mv(S/M)$ gives $\partial V/\partial S=v'(S/M)$ and $V/S=v(s)/s$.\par}

(b) Fix any policy and its share path $s_t$, which is nondecreasing, absolutely continuous, and on $[0,T]$ takes values in the compact set $[s_0,s(T)]\subset(0,\bar s)$. Since $\hat v\in C^1$, $J_t:=e^{-(r-G)t}\hat v(s_t)$ is absolutely continuous with $\dot J_t=e^{-(r-G)t}[\hat v'(s_t)\dot s_t-(r-G)\hat v(s_t)]$ almost everywhere. For every $\sigma\ge 0$, write $\Delta=\ell(s)\sigma^{\alpha}$ and $m=m_0-\sigma-k(G+\Delta)$; then the maximand in \eqref{eq:hjb} equals $m+\hat v'(s)\Delta$, so \eqref{eq:hjb} gives $(r-G)\hat v(s)\ge s\,m+\hat v'(s)\,s\Delta$. Along the path $s_t\Delta_t=\dot s_t$, hence $\dot J_t\le-e^{-(r-G)t}s_tm_t$ almost everywhere, and integrating,
\begin{equation}\label{eq:verif}
\int_0^Te^{-(r-G)t}s_tm_t\,dt\;\le\;\hat v(s_0)-e^{-(r-G)T}\hat v(s_T),
\end{equation}
where the left side may be $-\infty$. Because $\hat v$ is bounded and $r>G$, the last term tends to $0$, so the limit superior of the left side is at most $\hat v(s_0)$. Taking the supremum over policies gives $v\le\hat v$. Now suppose, as (H) requires, that the feedback generates a path $s_t$ from $s_0$. The map $s\mapsto\hat\sigma(s)=[\alpha\ell(s)(\hat v'(s)-k)_+]^{1/(1-\alpha)}$ is continuous, so $\sigma_t:=\hat\sigma(s_t)$ is continuous and locally bounded, hence admissible, and its induced share path is $s_t$ (the share equation has a unique solution for a given spend path). By (F2), $\hat\sigma(s)$ attains the maximum in \eqref{eq:hjb}, so the inequalities above hold with equality and \eqref{eq:verif} becomes an equality. Letting $T\to\infty$, the policy attains $\hat v(s_0)$. Hence $\hat v(s_0)\le v(s_0)$, so $\hat v(s_0)=v(s_0)$ and the policy is optimal. If this holds from every $s_0$, then $\hat v=v$.

(c) On the frontier, $g=G+\Delta$ and $m=m_0-\sigma-k(G+\Delta)$, so $R_X(g,m)-XG=X\Delta+m_0-\sigma-kG-k\Delta=m_0-\sigma-kG+(X-k)\Delta$, which is the maximand with $X=\hat v'(s)$. Maximizing it over $\sigma$ is therefore maximizing $R_X$ over the frontier at the current share, and by \autoref{prop:score}(a) also over the feasible set. Its maximum is $R_X^*(s)-XG$ by \autoref{prop:score}(e), for every $X$ when $\bar\sigma=\infty$. Under (H), $\hat v=v$ by (b), so $X^{\mathrm{dyn}}(s)=v'(s)=\partial V/\partial S$ by (a), and \eqref{eq:hjb} reads $(r-G)v(s)/s=R^*_{X^{\mathrm{dyn}}(s)}(s)-G\,X^{\mathrm{dyn}}(s)$.

(d) \emph{The root.} For $\upsilon\ge 0$ let $\Pi_\lambda(\upsilon):=C_\alpha[\lambda(\upsilon-k)_+]^{1/(1-\alpha)}$ and $\mathcal H(\upsilon):=(r-G)\upsilon-(m_0-kG)-\Pi_\lambda(\upsilon)$, so that \eqref{eq:gordon} is $\mathcal H(v_0)=0$. By the proof of \autoref{prop:envelope}, $\Pi_\lambda$ is continuously differentiable with $\Pi_\lambda'(\upsilon)=\Delta^*(\lambda,\upsilon)$. Since $k\ge 0$, $\mathcal H(0)=-(m_0-kG)<0$, so the smallest root $v_0$ is positive and $\mathcal H<0$ on $[0,v_0)$. Hence $\mathcal H'(v_0)\ge 0$, and simplicity gives $\mathcal H'(v_0)=(r-G)-\Delta^\circ>0$, that is, $g^\circ<r$. By \autoref{prop:score}(e), $\Pi_\lambda(v_0)=(v_0-k)_+\Delta^\circ-\sigma^\circ$, and since $\Delta^\circ=\sigma^\circ=0$ when $v_0\le k$, $\mathcal H(v_0)=0$ gives $(r-G)v_0=m_0-kG+(v_0-k)\Delta^\circ-\sigma^\circ$, that is, $(r-G-\Delta^\circ)v_0=m_0-kg^\circ-\sigma^\circ=m^\circ$. So $v_0=m^\circ/(r-g^\circ)$, and $m^\circ>0$.

\emph{Upper bound.} For every $s\in(0,\bar s)$ and $\sigma\ge 0$, $\ell(s)\le\lambda$ gives $(v_0-k)\ell(s)\sigma^{\alpha}\le(v_0-k)_+\lambda\sigma^{\alpha}$, so by (F2)
\[
m_0-\sigma-kG+(v_0-k)\ell(s)\sigma^{\alpha}\;\le\;m_0-kG+\Pi_\lambda(v_0)\;=\;(r-G)v_0 .
\]
Thus $\hat w(s):=v_0s$, which is bounded on $(0,\bar s)$, satisfies \eqref{eq:hjb} with ``$\ge$'' in place of ``$=$''. The first half of the proof of (b) uses only this inequality and boundedness, so $v(s)\le v_0s$.

\emph{Lower bound and limit.} If $v_0\le k$, then $v_0=(m_0-kG)/(r-G)$ and (a) gives $v(s)\ge v_0s$, so $v(s)=v_0s$. If $v_0>k$, use the constant policy $\sigma_t\equiv\sigma^\circ$. Its share path is logistic: $\hat s_t=\hat s_0e^{\Delta^\circ t}/(1-\hat s_0+\hat s_0e^{\Delta^\circ t})$ with $\hat s=s/\bar s$. Since $\Delta_t=\Delta^\circ h(s_t)\le\Delta^\circ$ and $k\ge 0$, $m_t=m_0-\sigma^\circ-kG-k\Delta_t\ge m^\circ>0$. Hence
\[
\frac{v(s_0)}{s_0}\;\ge\;m^\circ\int_0^\infty e^{-(r-G)t}\frac{s_t}{s_0}\,dt\;=\;m^\circ\int_0^\infty\frac{e^{-(r-G-\Delta^\circ)t}}{1+\hat s_0(e^{\Delta^\circ t}-1)}\,dt .
\]
As $s_0\downarrow 0$ the integrand increases pointwise to $e^{-(r-G-\Delta^\circ)t}$, which is integrable because $\Delta^\circ<r-G$. By monotone convergence the right side tends to $m^\circ/(r-G-\Delta^\circ)=v_0$. With the upper bound, $v(s)/s\to v_0$.

\emph{Marginal $q$.} Under (H), $\hat v=v$ by (b). By (F2), \eqref{eq:hjb} is equivalent to $(r-G)\hat v(s)/s-(m_0-kG)=C_\alpha[\ell(s)(\hat v'(s)-k)_+]^{1/(1-\alpha)}$. As $s\to 0$ the left side tends to $(r-G)v_0-(m_0-kG)=\Pi_\lambda(v_0)$. If $v_0>k$, this limit is positive, so for small $s$, $\hat v'(s)>k$ and $\ell(s)(\hat v'(s)-k)=[\,\cdot\,/C_\alpha]^{1-\alpha}\to\lambda(v_0-k)$; since $\ell(s)\to\lambda>0$, $\hat v'(s)\to v_0$. If $v_0\le k$, then $\hat v(s)=v(s)=v_0s$ and $\hat v'\equiv v_0$.

(e) Suppose $\mathcal H$ has no positive root. Since $\mathcal H(0)<0$ and $\mathcal H$ is continuous, $\mathcal H<0$ on $[0,\infty)$. The excess growth $\Delta^*(\lambda,\upsilon)$ is continuous in $\upsilon$, zero for $\upsilon\le k$ and unbounded as $\upsilon\to\infty$, so there is $\upsilon_1>k$ with $\Delta^*(\lambda,\upsilon_1)=r-G$. For $\upsilon\ge0$ write $g(\upsilon)=G+\Delta^*(\lambda,\upsilon)$ and $m(\upsilon)=m_0-kg(\upsilon)-\sigma^*(\lambda,\upsilon)$; as in (d), $\mathcal H(\upsilon)=(r-g(\upsilon))\upsilon-m(\upsilon)$. At $\upsilon_1$, $m(\upsilon_1)=-\mathcal H(\upsilon_1)>0$. By continuity, for $\upsilon<\upsilon_1$ close to $\upsilon_1$, $g(\upsilon)<r$ and $m(\upsilon)>m(\upsilon_1)/2$, so the Gordon multiple $m(\upsilon)/(r-g(\upsilon))$ of the constant policy $\sigma\equiv\sigma^*(\lambda,\upsilon)$ tends to infinity as $\upsilon\uparrow\upsilon_1$. The lower-bound argument of (d), which uses only $g(\upsilon)<r$ and $m(\upsilon)>0$, gives $\liminf_{s\to0}v(s)/s\ge m(\upsilon)/(r-g(\upsilon))$ for each such $\upsilon$. Hence $v(s)/s\to\infty$.

\emph{The HJB equation as an ODE.} Where the optimal spend is positive, (F2) shows that \eqref{eq:hjb} is equivalent to
\[
\hat v'(s)=k+\frac{1}{\ell(s)}\Bigl[\frac{(r-G)\hat v(s)/s-(m_0-kG)}{C_\alpha}\Bigr]^{1-\alpha},
\]
an explicit first-order equation that is singular at both ends of $(0,\bar s)$: at $s=0$, where $u:=\hat v/s$ satisfies $su'=\Phi(u,s)-u$ with $\Phi(u,s):=k+\ell(s)^{-1}\{[(r-G)u-(m_0-kG)]/C_\alpha\}^{1-\alpha}$, which has a singular point at $(0,v_0)$ with $\partial\Phi/\partial u=(r-G)/\Delta^\circ>1$ there, so every solution near it behaves like $u-v_0\propto s^{(r-G)/\Delta^\circ-1}$ and the boundary behavior $\hat v(s)/s\to v_0$ of (d) does not pin one down; and at $s=\bar s$, where $\ell(s)\to0$. A proof of (H) would have to select the solution through the terminal singular point and show that it is bounded and $C^1$ with $(r-G)\hat v(s)/s>m_0-kG$ on $(0,\bar s)$; we have not done this. \qed

\subsection*{Proof of \autoref{lem:valuation2}}

(a) Let $V=S_0[m\mathcal N_T(\eta)+m_{LR}\eta^T\mathcal M_\infty]$ with $m_{LR}=\theta m+(1-\theta)\bar m$. Then $\partial V/\partial m=S_0[\mathcal N_T(\eta)+\theta\eta^T\mathcal M_\infty]$. Since $\partial\eta/\partial g=1/(1+r)=\eta/(1+g)$,
\[
\frac{\partial V}{\partial g}=\frac{S_0}{1+g}\bigl[m\,\eta\mathcal N_T'(\eta)+m_{LR}T\eta^T\mathcal M_\infty\bigr],\qquad \eta\mathcal N_T'(\eta)=\sum_{t=1}^Tt\,\eta^t,
\]
and dividing gives \eqref{eq:X2}.

(b) With $\theta=1$ and $\mathcal M_\infty=0$, \eqref{eq:X2} is $m\sum_tt\eta^t/[(1+g)\sum_t\eta^t]=mD/(1+g)$. $D$ is the mean of $t$ under the weights $w_t\propto\eta^t$. Adding a date $T+1$ with positive weight raises the mean, because $T+1$ exceeds every existing date, so $D$ is strictly increasing in $T$. For fixed $T$, $dD/d\ln\eta=\operatorname{Var}_w(t)$, because the weights form an exponential family in $\ln\eta$; the variance is positive for $T\ge 2$ and zero for $T=1$. Since $\eta$ is strictly increasing in $g$ and strictly decreasing in $r$, the monotonicity claims follow. With $\theta=1$ and $\mathcal M_\infty>0$, $m_{LR}=m$ and \eqref{eq:X2} becomes $m[\sum_tt\eta^t+T\eta^T\mathcal M_\infty]/[(1+g)(\mathcal N_T+\eta^T\mathcal M_\infty)]=mD_{TV}/(1+g)$.

(c) If $m=0$, $\theta<1$ and $\bar m>0$, the numerator of \eqref{eq:X2} is $(1-\theta)\bar mT\eta^T\mathcal M_\infty>0$. For monotonicity in $\theta$, write the numerator as $N(\theta)=m\Sigma_1+[\bar m+\theta(m-\bar m)]T\mathcal M_\infty\eta^T$ with $\Sigma_1=\sum_tt\eta^t$, and the denominator as $(1+g)\mathcal D(\theta)$ with $\mathcal D(\theta)=\mathcal N_T+\theta\mathcal M_\infty\eta^T$. Then
\[
\frac{d}{d\theta}\frac{N}{\mathcal D}=\frac{(m-\bar m)T\mathcal M_\infty\eta^T\,\mathcal D-N\,\mathcal M_\infty\eta^T}{\mathcal D^2}.
\]
Let $0\le m\le\bar m$ and $\bar m>0$. Then $\mathcal D>0$, $\mathcal M_\infty\eta^T>0$ and $N\ge 0$, because $m\ge 0$ and $m_{LR}\ge\min\{m,\bar m\}\ge 0$. If $m<\bar m$, the first term of the numerator is strictly negative and the second is nonpositive. If $m=\bar m$, the first term is zero and $N=m\Sigma_1+mT\mathcal M_\infty\eta^T>0$, so the second term is strictly negative. In both cases the derivative is strictly negative.

(d) With $\theta=1$, $\mathcal M_\infty=0$ and $g<r$, $\eta\in(0,1)$, so $\sum_{t\ge 1}\eta^t=\eta/(1-\eta)$ and $\sum_{t\ge 1}t\eta^t=\eta/(1-\eta)^2$. Hence $D\to 1/(1-\eta)=(1+r)/(r-g)$ as $T\to\infty$, and $X^{\mathrm{scr}}\to m(1+r)/[(1+g)(r-g)]$. Also $V/S_0\to m\eta/(1-\eta)=m(1+g)/(r-g)$, and the ratio of the two limits is $(1+r)/(1+g)^2$.

(e) With $T=1$, $\mathcal N_1(\eta)=\sum_tt\eta^t=\eta$, and with $\theta=0$, $m_{LR}=\bar m$, so \eqref{eq:X2} is $[m\eta+\bar m\eta\mathcal M_\infty]/[(1+g)\eta]=(m+\bar m\mathcal M_\infty)/(1+g)$. Also $V/S_0=\eta(m+\bar m\mathcal M_\infty)$, and $\eta(1+g)=(1+g)^2/(1+r)$, which gives the second form. \qed

\subsection*{Proof of \autoref{prop:industry}}

(i) Every firm can choose $\sigma=0$, so $G\in\mathcal G_i$ for all $i$ and $\mathcal G=[G,\sup_i(G+\Delta_i^{\max})]$ (closed at the right end if the supremum is attained). Let $g_1<g_2$ in $\mathcal G$. Any firm $i$ with $g_2\in\mathcal G_i$ also has $g_1\in\mathcal G_i$, because $\mathcal G_i$ is an interval containing $G$, and $m_i^*(g_1)>m_i^*(g_2)$ by \autoref{prop:frontier}(c). Taking suprema over the (larger) set of firms feasible at $g_1$ gives $\Phi(g_1)\ge\Phi(g_2)$. If the supremum at $g_2$ is attained by some firm $j$, then $\Phi(g_1)\ge m_j^*(g_1)>m_j^*(g_2)=\Phi(g_2)$.

(ii) With common $k$ and $X>k$, \autoref{prop:score}(b) gives $g_\omega^c=G+\Delta^*(\ell_\omega,X)$, and $\Delta^*$ is strictly increasing in $\ell$, so $g_\omega^c$ is strictly increasing in $\omega$. The chosen margin is $m_\omega^c=m_{0\omega}-kG-[k+\alpha y]\Delta_\omega^*$ by \autoref{prop:score}(b), which gives \eqref{eq:locus}.

If $m_{0\omega}$ is nondecreasing, then for every $\omega$ and every $g\in\mathcal G_\omega$, $g\in\mathcal G_{\bar\omega}$ (because $\ell_{\bar\omega}\ge\ell_\omega$) and
\[
m_\omega^*(g)=m_{0\omega}-kg-\ell_\omega^{-\beta}(g-G)^{\beta}\le m_{0\bar\omega}-kg-\ell_{\bar\omega}^{-\beta}(g-G)^{\beta}=m_{\bar\omega}^*(g).
\]
Hence $\Phi=m_{\bar\omega}^*$ on $\mathcal G=\mathcal G_{\bar\omega}$, which is strictly decreasing.

Under \eqref{eq:quality}, \eqref{eq:locus} gives $m_{\omega'}^c\ge m_\omega^c$ for $\omega<\omega'$, so the locus is nondecreasing in $g$. For any $g$ in the domain of $\Psi$, the top firm satisfies $g_{\bar\omega}^c=\sup_\omega g_\omega^c\ge g$, so it belongs to $\{\omega:g_\omega^c\ge g\}$, and its margin is the largest in that set. Thus $\Psi(g)=m_{\bar\omega}^c$.

For the quantile frontiers, let $\Omega=[\underline\omega,\bar\omega]$ carry a continuous, strictly increasing distribution function $F_\Omega$. The map $\omega\mapsto g_\omega^c$ is continuous and strictly increasing, so for $g\in[g_{\underline\omega}^c,g_{\bar\omega}^c)$ the set $\{\omega:g_\omega^c\ge g\}$ is an interval $[\omega(g),\bar\omega]$, with $\omega(g)$ continuous and strictly increasing in $g$. Under strict \eqref{eq:quality}, $\omega\mapsto m_\omega^c$ is continuous and strictly increasing, so the $\chi$-quantile of $m^c$ on $[\omega(g),\bar\omega]$ is $m^c$ evaluated at the $\chi$-quantile of $\omega$ conditional on $\omega\ge\omega(g)$, namely $\omega_\chi(g)=F_\Omega^{-1}\bigl(\chi+(1-\chi)F_\Omega(\omega(g))\bigr)$. This is strictly increasing in $\omega(g)$ for $\chi<1$, because $F_\Omega$ is strictly increasing. Hence $\Psi_\chi$ is strictly increasing on $[g_{\underline\omega}^c,g_{\bar\omega}^c)$.

(iii) For the first equality, $\sup_{g\in\mathcal G}[Xg+\Phi(g)]=\sup_g\sup_{i:g\in\mathcal G_i}[Xg+m_i^*(g)]=\sup_i\sup_{g\in\mathcal G_i}[Xg+m_i^*(g)]=\sup_iR_i^*$, by interchanging suprema and using \autoref{prop:score}(a) (or \autoref{rem:cap}) for each firm. For the second, each firm chooses its optimum, so $R_i^c=R_i^*$. For any $g$ and any firm $i$ with $g_i^c\ge g$, $Xg+m_i^c\le Xg_i^c+m_i^c=R_i^*$ because $X\ge 0$; taking suprema gives $\sup_g[Xg+\Psi(g)]\le\sup_iR_i^*$. Conversely, evaluating at $g=g_i^c$ gives $Xg_i^c+\Psi(g_i^c)\ge Xg_i^c+m_i^c=R_i^*$ for every $i$. Under the assumptions of (ii), $\Psi\equiv m_{\bar\omega}^c$ is constant and $X>k\ge 0$, so $Xg+\Psi(g)$ is strictly increasing in $g$ and is maximized only at the right end of its domain, $g=g_{\bar\omega}^c$, where it equals $R_{\bar\omega}^*$. At that point $\Psi$ is flat while $m_{\bar\omega}^{*\prime}(g_{\bar\omega}^c)=-X$, so the maximum is not a tangency of an iso-score line with $\Psi$.

(iv) and (v) are the summaries of \autoref{prop:industry-full}(iv) and (v), proved next; for (v), \eqref{eq:bW} can be written $\mathcal W=w_\ell\cdot\{-[k+\alpha(X-k)]\}+w_{\tilde g}\cdot\mathcal L+C_\ell/(V_\ell+V_{\tilde g})$ with weights $w_\ell=V_\ell/(V_\ell+V_{\tilde g})$ and $w_{\tilde g}=V_{\tilde g}/(V_\ell+V_{\tilde g})$. \qed

\subsection*{Proof of \autoref{prop:industry-full}}

(iv) By \autoref{prop:score}(b), firm $\ell$ chooses $g^c=G+\Delta^*(\ell,X)$ and $m^c=m_0(\ell)-kG-(k+\alpha y)\Delta^*(\ell,X)=\Upsilon-kG+\xi\Delta^*(\ell,X)$, so its chosen point lies on the line $m^c=\Upsilon-kG+\xi(g^c-G)$. Since $\ell\mapsto\Delta^*(\ell,X)$ is a continuous increasing bijection from $[\ell_{lo},\ell_{hi}]$ onto its image, the chosen points fill the segment with $g^c\in[g_{lo},g_{hi}]$. On $(g_{lo},g_{hi})$, $\Psi(g)$ is the supremum of $\Upsilon-kG+\xi(g'-G)$ over $g'\in[g,g_{hi}]$, which equals $\Upsilon-kG+\xi(g_{hi}-G)$ if $\xi\ge 0$ and $\Upsilon-kG+\xi(g-G)$ if $\xi<0$; its slope is $\min\{\xi,0\}$. At the chosen point, $m_i^{*\prime}(g_i^c)=-X$ by tangency, and $m_i^{*\prime\prime}(g_i^c)=-\beta(\beta-1)\kappa_i(\Delta_i^*)^{\beta-2}=-(\beta-1)\cdot\beta\kappa_i(\Delta_i^*)^{\beta-1}/\Delta_i^*=-(\beta-1)y/\Delta_i^*=-(1-\alpha)y/(\alpha\Delta_i^*)$, using the first-order condition $\beta\kappa_i(\Delta_i^*)^{\beta-1}=y$.

For the second claim, fix $\alpha'\in(0,1)$ and $X'>k$, let $y'=X'-k$, and for each excess growth $\Delta\in[g_{lo}-G,g_{hi}-G]$ define the firm with opportunity $\ell'(\Delta)=\Delta^{1-\alpha'}(\alpha'y')^{-\alpha'}$ and intercept $m_0'(\Delta)=\Upsilon+[\xi+k+\alpha'y']\Delta$. Inverting \eqref{eq:Dstar} shows that this firm, maximizing $R_{X'}$ with elasticity $\alpha'$, chooses $\Delta^*=\Delta$ with spend $\alpha'y'\Delta\le y'(g_{hi}-G)\le\bar\sigma$, so its cap is slack. Its chosen margin is $\Upsilon-kG+\xi\Delta$ by \autoref{prop:score}(b), so the chosen points again fill the same segment, while the slope and curvature at the chosen points are $-X'$ and $-(1-\alpha')y'/(\alpha'\Delta)$ by the first paragraph.

(v) Write $\Delta^*(\cdot)$ and $\sigma^*(\cdot)$ for $\Delta^*(\cdot,X)$ and $\sigma^*(\cdot,X)$. With $\ell_t=\bar\ell+\varepsilon z_{2t}$, the planned growth is $g_t^e=G+\Delta^*(\bar\ell+\varepsilon z_{2t})$ and $g_t=g_t^e+\varepsilon z_{3t}$. Because $\Delta^*$ and $\sigma^*$ are smooth on a neighborhood of $\bar\ell>0$ and $z$ is bounded, Taylor's theorem gives, uniformly in $z$,
\[
g_t-\bar g^e=\varepsilon\bigl(\Delta^{*\prime}z_{2t}+z_{3t}\bigr)+O(\varepsilon^2),
\]
with derivatives evaluated at $\bar\ell$. For the margin, define
\[
\mathcal M(m_0,\ell,\tilde g)=m_0+a^f-\frac{(a^f+\sigma^*(\ell))(1+g^e(\ell))}{1+g^e(\ell)+\tilde g}-k\bigl(g^e(\ell)+\tilde g\bigr).
\]
 Its partial derivatives at $(\bar m_0,\bar\ell,0)$ are $\partial\mathcal M/\partial m_0=1$; $\partial\mathcal M/\partial\tilde g=(a^f+\sigma^*(\bar\ell))/(1+\bar g^e)-k=\mathcal L$; and, at $\tilde g=0$, $\mathcal M=m_0-\sigma^*(\ell)-kg^e(\ell)$, so $\partial\mathcal M/\partial\ell=-\sigma^{*\prime}-k\Delta^{*\prime}=-(k+\alpha y)\Delta^{*\prime}$, using $\sigma^*=\alpha y\Delta^*$. Hence $m_t-\bar m=\varepsilon(z_{1t}-(k+\alpha y)\Delta^{*\prime}z_{2t}+\mathcal Lz_{3t})+O(\varepsilon^2)$ uniformly. Let $\hat g:=\Delta^{*\prime}z_2+z_3$ and $\hat m:=z_1-(k+\alpha y)\Delta^{*\prime}z_2+\mathcal Lz_3$. The population covariance and variance satisfy $\operatorname{Cov}(m_t,g_t)=\varepsilon^2\operatorname{Cov}(\hat m,\hat g)+O(\varepsilon^3)$ and $\operatorname{Var}(g_t)=\varepsilon^2\operatorname{Var}(\hat g)+O(\varepsilon^3)$, because all remainders are uniformly bounded. Since $z_3$ is uncorrelated with $(z_1,z_2)$,
\[
\operatorname{Cov}(\hat m,\hat g)=C_\ell-(k+\alpha y)V_\ell+\mathcal LV_{\tilde g},\qquad \operatorname{Var}(\hat g)=V_\ell+V_{\tilde g}.
\]
Dividing and letting $\varepsilon\to 0$ gives \eqref{eq:bW}, because $V_\ell+V_{\tilde g}>0$. If only opportunity shocks are present ($V_{\tilde g}=0$ and $z_1\equiv 0$), $\mathcal W=-(k+\alpha y)$. The sign condition follows because the denominator is positive. In the numerical example after the proposition, $\sigma^*=\alpha y\Delta^*=0.03X$, $1+\bar g^e=1.2$, and $\mathcal W=0$ with $C_\ell=0$ requires $V_{\tilde g}/(V_\ell+V_{\tilde g})=0.3X/(0.3X+\mathcal L)$, which is $0.75/1.0625\approx0.71$, $1.5/1.875=0.80$ and $2.7/3.175\approx0.85$ at $X=2.5$, 5 and 9. \qed

\subsection*{Proof of \autoref{rem:segments}}

Firm 1 is feasible at $g=0.1$ (given $\ell_1\bar\sigma^{\alpha}\ge 0.1$) with $m_1^*(0.1)=0.2-\kappa_1(0.1)^{\beta}<0.2$. Firm 2 is not feasible at $g=0.1<G_2$, because its growth is at least $G_2$. Hence $\Phi(0.1)=m_1^*(0.1)<0.2$. At $g=0.3$, firm 1 attains at most $m_1^*(0.3)<0.2$ (if feasible), and firm 2 attains $m_2^*(0.3)=0.5$, so $\Phi(0.3)=0.5>\Phi(0.1)$. \qed

\subsection*{Proof of \autoref{prop:entry}}

(a) Fix $n$ firms. Owning capital gives the same present value as renting it along any nondecreasing output path $x(t)$ with $e^{-rt}x(t)\to0$: the owner invests $k_ux(0)$ at entry and $k_u\,dx(t)$ thereafter, and integration by parts gives $k_ux(0)+\int_{0+}^\infty e^{-rt}k_u\,dx(t)=\int_0^\infty re^{-rt}k_ux(t)\,dt$, the present value of rental payments. Monotonicity matters because capital, once bought, is not sold. The same argument applies to customers, acquired at $a_c$ per unit of output when output grows, against the flow cost $ra_c$. We therefore work with the flow form, in which a firm's flow profit at date $t$ is $(A-B(x+Q_-)e^{-Gt}-\hat c)x$ with $\hat c=u+r(k_u+a_c)$ and $Q_-$ the others' output. The best response is $x=[(A-\hat c)e^{Gt}/B-Q_-]/2$, and the symmetric equilibrium is $x_n(t)=(A-\hat c)e^{Gt}/[B(n+1)]$, with price $P_n=A-Bnx_ne^{-Gt}=(A+n\hat c)/(n+1)$, constant over time, so output grows at $G$ and is nondecreasing. Flow profit is $(P_n-\hat c)x_n=\pi_ne^{Gt}$ with $\pi_n=(A-\hat c)^2/[B(n+1)^2]$, and an entrant's net present value is $\int_0^\infty e^{-rt}\pi_ne^{Gt}dt-F=\pi_n/(r-G)-F$. The Lerner index is $(P_n-\hat c)/P_n=[(A-\hat c)/(n+1)]/[(A+n\hat c)/(n+1)]=(A-\hat c)/(A+n\hat c)$. For a firm that owns its capital and acquires its customers, revenue is $P_nx$, variable cost is $ux$, acquisition spend is $a_c\dot x=a_cGx$ (an expense), and the capital stock is $k_ux$, so $k_n=k_u/P_n$ and $k_{a,n}=a_c/P_n$. The operating margin is $(P_n-u-a_cG)/P_n=(P_n-\hat c)/P_n+rk_u/P_n+(r-G)a_c/P_n=L_n+rk_n+(r-G)k_{a,n}$. Investment is $k_u\dot x=k_uGx$, so the FCF margin is $m_{op}-Gk_n=(r-G)(k_n+k_{a,n})+L_n$. Since revenue grows at $G$, the score is $XG+(r-G)(k_n+k_{a,n})+L_n$.

(b) An entrant enters if and only if $\pi_n/(r-G)\ge F$, that is, $(A-\hat c)^2/[B(n+1)^2]\ge\hat F$, or $n+1\le N_F$. Because $\pi_n$ is decreasing in $n$, free entry gives the largest $n$ with $n+1\le N_F$, namely $n^*=\lfloor N_F\rfloor-1\ge 1$. Positivity of $L_{n^*}$ holds because $A>\hat c$. For the upper bound, $A+n^*\hat c>\hat c+n^*\hat c=\hat c\lfloor N_F\rfloor>\hat c(N_F-1)>0$, because $A>\hat c>0$ and $\lfloor N_F\rfloor>N_F-1\ge 1$; hence $L_{n^*}=(A-\hat c)/(A+n^*\hat c)<(A-\hat c)/[\hat c(N_F-1)]$. As $F\to 0$, $\hat F\to 0$, so $N_F\to\infty$, $n^*\to\infty$, $L_{n^*}\to 0$ and $P_{n^*}\to\hat c$; hence $k_{n^*}+k_{a,n^*}\to(k_u+a_c)/\hat c$ and the score tends to $XG+(r-G)(k_u+a_c)/\hat c$. For $X=1$, $G+(r-G)\hat k=G(1-\hat k)+r\hat k$.

(c) With the input price $w$ per unit, the full unit cost is $\hat c+w$, and the Cournot quantities from (a) give total output $Q_n(w)=n(A-\hat c-w)e^{Gt}/[B(n+1)]$ whenever $w<A-\hat c$. The input market clears with $w>0$ if and only if $Q_n(0)>Ze^{Gt}$, that is, $n(A-\hat c)/[B(n+1)]>Z$. Since $n/(n+1)\ge 1/2$ and $Z<(A-\hat c)/(2B)$, this holds for every $n\ge 1$. Solving $Q_n(w)=Ze^{Gt}$ gives $w_n=A-\hat c-BZ(n+1)/n>0$, and the price is $P=A-BZ$, independent of $n$. Each firm produces $x=Ze^{Gt}/n$ and earns flow profit net of all costs (including the input at price $w_n$) $(P-\hat c-w_n)x=[BZ/n]\cdot Ze^{Gt}/n=BZ^2e^{Gt}/n^2$; the net Lerner index is $L_n^{net}=(P-\hat c-w_n)/P=BZ/(nP)$. A firm that owns its input requirement pays nothing for the input, so, as in (a), its operating margin is $(P-u-a_cG)/P=rk_u/P+(r-G)a_c/P+w_n/P+(P-\hat c-w_n)/P=rk+(r-G)k_a+\rho_n+L_n^{net}$, with $k=k_u/P$, $k_a=a_c/P$ and $\rho_n=w_n/P$. A firm that buys the input has $m_{op}=rk+(r-G)k_a+L_n^{net}$. An entrant buys the input, so entry is profitable if and only if $BZ^2/[n^2(r-G)]\ge F$, that is, $n\le Z\sqrt{B/\hat F}$. Free entry therefore gives $n^*=\lfloor Z\sqrt{B/\hat F}\rfloor$, which is at least $1$ under the stated assumption and tends to infinity as $F\to 0$. Then $w_{n^*}\to A-\hat c-BZ$, $\rho_{n^*}\to(A-BZ-\hat c)/(A-BZ)=1-\hat c/(A-BZ)=\rho$, and $L_{n^*}^{net}\to 0$. FCF margins subtract $Gk$ as in (a), which gives \eqref{eq:Rrent}.

(d) By (c), $\rho=1-\hat c/P=1-u/P-r(k+k_a)$ at the limit price $P=A-BZ$. If $u/P\ge c_c+f+a_{\min}+c_{\mathrm{full}}$, then $\rho\le1-c_c-f-a_{\min}-c_{\mathrm{full}}-r(k+k_a)$, and the input owner's limit score satisfies $XG+(r-G)(k+k_a)+\rho\le1-c_c-f-a_{\min}-c_{\mathrm{full}}+XG-G(k+k_a)$, which is $R_{\max}$ of \autoref{cor:ceiling} with $\tau=0$, $\delta_K=0$ and capital $k+k_a$. \qed

\subsection*{Proof of \autoref{rem:bertrand}}

Under homogeneous Bertrand competition with identical full unit cost $\hat c$, two or more active firms price at $\hat c$ and earn zero flow profit, so the net value of entry for a second firm is $-F<0$. Hence at most one firm enters, and a single entrant charges the monopoly price, which it does provided monopoly profit covers $F$. Under contestability with a flow fixed cost $\hat F$ that is not sunk, any price above average cost $\hat c+\hat F/x$ invites hit-and-run entry, so in equilibrium $P=\hat c+\hat F/x$ and $L=(P-\hat c)/P=\hat F/(Px)\to 0$ as $F\to 0$. A firm that rents its capital and pays for its customers as a flow holds neither stock, so $k=k_a=0$ and the score in \autoref{prop:entry}(a) is $XG+L_n$, with limit $XG$. Under monopolistic competition with differentiated varieties, each firm's markup over marginal cost is set by the elasticity of substitution between varieties, and free entry drives profit to zero by adding varieties until the markup just covers the fixed cost, so the markup does not vanish as $F\to0$ unless varieties become perfect substitutes. \qed

\subsection*{Proof of \autoref{lem:prices}}

(a) $\gamma=1-f-u/p(u)$. Differentiating, $d\gamma/du=-[p-up']/p^2=-(1/p)[1-up'/p]=-(1-\epsilon_p)/p$. A reduction in $u$ raises $\gamma$ if and only if $d\gamma/du<0$, that is, $\epsilon_p<1$. Cournot with linear demand: from the proof of \autoref{prop:entry}(a), $p=(A+n(u+rk_u))/(n+1)$, so $p'=n/(n+1)$ and $\epsilon_p=\tfrac{n}{n+1}u/p<1$. Monopoly with constant elasticity $\eta_D>1$: $p=\tfrac{\eta_D}{\eta_D-1}(u+rk_u)$, so $\epsilon_p=u/(u+rk_u)$. Perfect competition: $p=u+rk_u$, so $\epsilon_p=u/(u+rk_u)$. In both of the last two cases $\epsilon_p\le 1$, with equality if and only if $rk_u=0$, that is, $k_u=0$ (given $r>0$). Finally $dp/du=(d\ln p/d\ln u)(p/u)=\epsilon_p p/u$; under constant-elasticity monopoly this is $\eta_D/(\eta_D-1)>1$.

(b) Nominal market size is $M_t=p_tQ_t$ with $Q_t=Q_0e^{G_Qt}p_t^{-\eta_D}$ and $p_t=p_0e^{-\delta_pt}$. Then $\ln M_t=\text{const}+G_Qt-(1-\eta_D)\delta_pt$, so $G=d\ln M/dt=G_Q-(1-\eta_D)\delta_p$. The sign claims are immediate. \qed

\subsection*{Proof of \autoref{prop:ai}}

By \autoref{prop:envelope}, $R_X^*$ is continuously differentiable in $(m_0,G,k,\lambda)$, so the chain rule gives $dR_X^*/d\zeta=m_0'+(X-k)G'-g^*k'+(X-k)(\Delta^*/\lambda)\lambda'$, with $m_0'=\gamma'-a'$. If the parameters are continuously differentiable, so is $\zeta\mapsto R_X^*$, and a positive (negative) derivative at $\zeta$ persists on a neighborhood, on which $R_X^*$ is strictly increasing (decreasing). The expression for $\gamma'$ follows from \autoref{lem:prices}(a) and the chain rule. \qed

\subsection*{Proof of \autoref{prop:ai-shift}}

(a) Denote old and new parameters by subscripts 0 and 1, with $\lambda_1\ge\lambda_0$, $G_1\ge G_0$, $k_1\le k_0$ and $m_{01}-k_1G_1\ge m_{00}-k_0G_0$. Since $h(s)$ is unchanged, $\ell_1\ge\ell_0$, so $\kappa_1\le\kappa_0$ and the new growth bound satisfies $\bar g_1-G_1\ge\bar g_0-G_0$. Take a feasible old point $(g,m)$, so $m\le m_0^*(g)$ with $g\in[G_0,\bar g_0]$. Case 1: $g\ge G_1$. Then $g-G_1\le g-G_0\le\bar g_0-G_0\le\bar g_1-G_1$, so $g$ is feasible under the new parameters. Write $m_j^*(g)=(m_{0j}-k_jG_j)-k_j(g-G_j)-\kappa_j(g-G_j)^{\beta}$ for $j=0,1$. Then $m_1^*(g)-m_0^*(g)$ is the sum of three nonnegative terms: $(m_{01}-k_1G_1)-(m_{00}-k_0G_0)\ge 0$; $k_0(g-G_0)-k_1(g-G_1)\ge 0$, because $k_0\ge k_1\ge 0$ and $g-G_0\ge g-G_1\ge 0$; and $\kappa_0(g-G_0)^{\beta}-\kappa_1(g-G_1)^{\beta}\ge 0$, for the same reason. Hence $(g,m_1^*(g))$ is new-feasible and dominates or equals $(g,m)$. Case 2: $g<G_1$. Then the new point $(G_1,m_1^*(G_1))$ has $G_1>g$ and $m_1^*(G_1)=m_{01}-k_1G_1\ge m_{00}-k_0G_0=m_0^*(G_0)\ge m_0^*(g)\ge m$, using \autoref{prop:frontier}(c). So $(g,m)$ is strictly dominated. In both cases every old feasible point is weakly dominated by a new feasible point. Since $X\ge 0$, a dominating point has a weakly higher score, so $R_{X,1}^*\ge R_{X,0}^*$. Reversing the roles of the two parameter vectors proves the reverse statement.

(b) The difference between the new and old frontier at $g$ is $dm_0-dk\,g$, which is positive for $g<dm_0/dk$ and negative for $g>dm_0/dk$. \qed

\subsection*{Proof of \autoref{rem:taxes}}

Write $m_0=m_0^{op}-\delta_Kk$, where $m_0^{op}$ is the margin intercept before the replacement spending $\delta_Kk$ that \autoref{ass:capital} includes in $a$. Taxable profit per dollar of revenue is $m_0^{op}-\sigma$ less the deductions for capital spending, growth and replacement alike, whose present value is $\psi$ per dollar. With linear taxes and full loss offset the after-tax FCF margin is $(1-\tau)(m_0^{op}-\sigma)-(1-\psi\tau)\delta_Kk-(1-\psi\tau)kg=(1-\tau)(m_0-\sigma)-(1-\psi)\tau\delta_Kk-(1-\psi\tau)kg$, where $\psi\tau kg$ is the present value of the tax shield on new capital. Let $\tilde\sigma=(1-\tau)\sigma$ be after-tax spend. Then $g=G+\ell\sigma^{\alpha}=G+\ell(1-\tau)^{-\alpha}\tilde\sigma^{\alpha}$, so the growth technology in terms of $\tilde\sigma$ has demand pull $\tilde\lambda=(1-\tau)^{-\alpha}\lambda$, and $\tilde\sigma\in[0,(1-\tau)\bar\sigma]$. The margin is $\tilde m_0-\tilde\sigma-\tilde kg$ with $\tilde m_0=(1-\tau)m_0-(1-\psi)\tau\delta_Kk$ and $\tilde k=(1-\psi\tau)k$. This is Assumptions~\ref{ass:growth}--\ref{ass:capital} with the substituted parameters, so every result derived from those assumptions applies. Explicitly, the taxed frontier is $\tilde m_0-\tilde kg-(\tilde\lambda h)^{-\beta}(g-G)^{\beta}$, and $(\tilde\lambda h)^{-\beta}=(1-\tau)\kappa$. For $\alpha=0.3$, $(1-\tau)^{-\alpha}$ equals $1.050$ at $\tau=0.15$ and $1.073$ at $\tau=0.21$. For \eqref{eq:psiT}, a dollar deducted ratably over $T$ years from the middle of the first tax year yields a half-year deduction of $1/(2T)$ in the first year, $1/T$ in each of the next $T-1$ years and $1/(2T)$ in the year after; discounting each year's deduction to the first, $\psi_T=1/(2T)+(1/T)\sum_{t=1}^{T-1}(1+r)^{-t}+(1+r)^{-T}/(2T)$, and the geometric sum is $[1-(1+r)^{-(T-1)}]/r$. At $r=0.10$, $\psi_5=0.796$ and $\psi_{15}=0.532$. \qed

\subsection*{Proof of \autoref{cor:ceiling}}

(a) By \autoref{rem:taxes} and \eqref{eq:Rinf}, the saturated score is $\tilde m_0+(X-\tilde k)G$ with $\tilde m_0=(1-\tau)m_0-(1-\psi)\tau\delta_Kk$ and $\tilde k=(1-\psi\tau)k$. With non-discretionary cost $a'=a'_{op}+\delta_Kk$, $m_0=1-c_h'-c_c'-f'-a'_{op}-\delta_Kk$, so $\tilde m_0=(1-\tau)(1-c_h'-c_c'-f'-a'_{op})-(1-\psi\tau)\delta_Kk$, which gives the expression for $R_\infty$. $R_{\max}$ is the same expression with total cost share $c_c+f+a_{\min}+c_{\mathrm{full}}$ in place of $C':=c_h'+c_c'+f'+a'_{op}$. Because $1-\tau>0$, $R_\infty\le R_{\max}$ if and only if $C'\ge c_c+f+a_{\min}+c_{\mathrm{full}}$, with equality of the scores if and only if the totals are equal. A firm at the full-automation structure has $c_c'=c_c$, $f'=f$ and $a'_{op}\ge a_{\min}+c_{\mathrm{full}}$, so $C'\ge c_c+f+a_{\min}+c_{\mathrm{full}}$, with equality if and only if $c_h'=0$ and $a'_{op}=a_{\min}+c_{\mathrm{full}}$. A firm with $C'$ below that total has $R_\infty>R_{\max}$.

(b) With $X=1$, $\psi=0$ and $\delta_K=0$, $R_{\max}$ reduces to \eqref{eq:Rmax}; at $k=0$ it is at least $\bar R$ if and only if $G\ge\bar R-(1-\tau)(1-c_c-f-a_{\min}-c_{\mathrm{full}})$. In general, $\partial R_{\max}/\partial\psi=\tau kG+\tau\delta_Kk$, so setting $\psi=1$ instead of $0$ adds $\tau kG$ when $\delta_K=0$, and a positive $\delta_K$ subtracts $(1-\psi\tau)\delta_Kk$.

(c) At saturation, $\Delta^*=0$ and $g^*=G$, and the saturated score, $R_\infty=m_0+(X-k)G$ as a function of the parameters $\vartheta=(m_0,k,\ln\lambda,G)$, has gradient $(1,-G,0,X-k)$; so a process is $X$-improving at saturation if and only if $\delta m_{0p}-G\delta k_p+(X-k)\delta G_p\ge0$, which is \autoref{def:attainable}(c) evaluated at $s=\bar s$, where $\Delta^*=0$ and $g^*=G$; at a share with headroom the two conditions differ when a process changes $\lambda$ or $k$. Because its directional derivative along each $\delta\vartheta_p$ is nonnegative at every point of $\mathcal B:=\{\vartheta_0+\sum_px_p\delta\vartheta_p:x\in[0,1]^N\}$, moving along $\delta\vartheta_p$ inside $\mathcal B$ never lowers it, so adding a process never lowers it, so the maximum over $\mathcal Z\subseteq\mathcal P$ is attained at $\mathcal P$. If some process has a negative directional derivative, dropping it raises the saturated score, so a partial program can do better.

The statement about $R_X^*(s)-R_\infty$ is \autoref{prop:saturation}(a)--(b) applied to the substituted parameters, and the integral bound is \autoref{prop:budget}(a), whose proof uses only $m\le\tilde m_0-\tilde kg$, $g\ge G$ and $\dot s/s=g-G$. The entries of Tables~\ref{tab:ceiling} and \ref{tab:fullcost} are direct evaluations of these formulas. \qed

\subsection*{Proof of \autoref{prop:amin}}

(a) A misstatement goes undetected if and only if every layer of the design fails. The cross-type common cause, with probability $p^{\mathrm{cc}}$, defeats every layer. Otherwise, for a type $j$ with $n_{oj}\ge1$ layers, all of them fail either through the common mode, with probability $p^{\mathrm{cm}}_j$, or, if it does not occur, each independently, with probability $p_j^{n_{oj}}$, so all fail with probability $p^{\mathrm{cm}}_j+(1-p^{\mathrm{cm}}_j)p_j^{n_{oj}}$; types fail independently, which gives the product. For one type the factor is strictly decreasing in $n_{oj}$, because $p_j<1$ and $p^{\mathrm{cm}}_j<1$, and it tends to $p^{\mathrm{cm}}_j$; since $p^{\mathrm{cc}}<1$, $P_o$ is strictly decreasing too. Each factor is at least $p^{\mathrm{cm}}_j$, which gives the lower bound, and the factors are strictly above $p^{\mathrm{cm}}_j$ at every finite depth. At given layers, $\partial P_o/\partial p^{\mathrm{cc}}=1-\prod_j[\cdot]\ge0$ and $\partial P_o/\partial p^{\mathrm{cm}}_j=(1-p^{\mathrm{cc}})(1-p_j^{n_{oj}})\prod_{i\ne j}[\cdot]\ge0$, and both probabilities are nonincreasing in $\theta$, so $P_o$ is nonincreasing in $\theta$.

(b) $\pi_o(D)\le\bar\pi$ if and only if $(1-P_o(D))^{N_o}\ge1-\bar\pi$, that is, $P_o(D)\le\bar P_o$. For a design that uses only types in $J$, (a) gives $P_o(D)>p^{\mathrm{cc}}(\theta)+(1-p^{\mathrm{cc}}(\theta))\prod_{j\in J}p^{\mathrm{cm}}_j(\theta)$ at its quality $\theta$, strictly because every factor exceeds its common-mode probability at finite depth and $p^{\mathrm{cc}}<1$. If the right side is at least $\bar P_o$ for every $\theta$, then $P_o(D)>\bar P_o$ and $D$ is not admissible. The statement about a single AI provider is the case $J=\{j\}$; since the layers' cost and number do not enter the bound, only $\theta$ or the addition of another type can lower it.

(c) The floor is the definition. The minimum exists when some design $D_0$ is admissible: only designs that cost no more than $D_0$ matter, and since $c_j>0$ and $\Theta$ is finite there are finitely many of them. Material misstatements enter only through the constraint, so an inadmissible design is excluded whatever its cost.

(d) The admissibility condition depends on $N_o$, $\bar\pi$ and the probabilities, none of which depends on $S$. Let $F_D$ be the dollar cost of design $D$ other than the expected cost of errors below materiality; the floor of $D$ is $F_D/S+\sum_o\omega_oL^{\mathrm{nm}}_oP_o(D)$, nonincreasing in $S$, so the minimum over the fixed admissible set is nonincreasing in $S$. For the limit, let $\bar\theta:=\max\Theta$, which minimizes the bracket because both probabilities are nonincreasing in $\theta$. For every admissible $D$, $P_o(D)\ge p^{\mathrm{cc}}(\bar\theta)+(1-p^{\mathrm{cc}}(\bar\theta))\prod_{j\in\mathcal J}p^{\mathrm{cm}}_j(\bar\theta)$, because unused types contribute a factor one and $p^{\mathrm{cm}}_j\le1$; so the floor is at least $\sum_o\omega_oL^{\mathrm{nm}}_o$ times the bracket at $\bar\theta$. Conversely, starting from an admissible design, raise $\theta$ to $\bar\theta$ and add layers of every type in $\mathcal J$; by (a) neither step raises any $P_o$, so the design stays admissible, and with enough layers $P_o$ is within any $\epsilon>0$ of the bracket for every $o$. For that fixed design $F_D/S\to0$, so the limit superior of the floor is at most the bound plus $\epsilon\sum_o\omega_oL^{\mathrm{nm}}_o$. For a single AI provider the bracket is $p^{\mathrm{cc}}(\bar\theta)+(1-p^{\mathrm{cc}}(\bar\theta))p^{\mathrm{cm}}_j(\bar\theta)>0$. \qed

\subsection*{Proof of \autoref{lem:fullcost}}

(a) Upkeep costs $\delta_Mc_{\rm bld}$ a year. Revalidations of $\varrho_{\rm rv}c_{\rm bld}$ every $T_{\rm up}$ years have the undiscounted annual equivalent $\varrho_{\rm rv}c_{\rm bld}/T_{\rm up}$. With discounting, payments $C$ at $T_{\rm up},2T_{\rm up},\dots$ have present value $C/(e^{rT_{\rm up}}-1)$, which equals that of the flow $rC/(e^{rT_{\rm up}}-1)$.

(b) Let $\mathcal C(c)=c+e^{-c/c_{\rm eff}}c^{\rm inc}$ on $c\ge0$. Then $\mathcal C''>0$ and $\mathcal C'(c)=1-(c^{\rm inc}/c_{\rm eff})e^{-c/c_{\rm eff}}$. If $c^{\rm inc}\le c_{\rm eff}$, $\mathcal C'(0)\ge0$, so the minimum is $\mathcal C(0)=c^{\rm inc}$. Otherwise $\mathcal C'(c^*)=0$ at $c^*=c_{\rm eff}\ln(c^{\rm inc}/c_{\rm eff})>0$, where the residual loss is $e^{-c^*/c_{\rm eff}}c^{\rm inc}=c_{\rm eff}$ and $\mathcal C(c^*)=c_{\rm eff}[1+\ln(c^{\rm inc}/c_{\rm eff})]$. Since $\ln x\le x/e$ for $x>0$, $c^*\le c_{\rm eff}(c^{\rm inc}/c_{\rm eff})/e=c^{\rm inc}/e$. If today's spending is optimal and positive, $c^{\rm sp}_0=c_{\rm eff}L_{\rm sec}$, so $c_{\rm eff}=c^{\rm sp}_0/L_{\rm sec}$. At full automation $c^{\rm inc}=\xi_{\rm exp}c^{\rm inc}_0=\xi_{\rm exp}c_{\rm eff}e^{L_{\rm sec}}>c_{\rm eff}$, so $c^{\rm risk}=c_{\rm eff}(1+L_{\rm sec}+\ln\xi_{\rm exp})$.

(c) An error escapes $n_r$ review layers on one provider if the review fails with the model, with probability $\chi$, or, if not, if each layer misses it independently, with probability $(1-d)^{n_r}$; so an ordinary error escapes with probability $\chi+(1-\chi)(1-d^{\rm o})^{n_r}$ and a fabrication with $\chi+(1-\chi)(1-d^{\rm h})^{n_r}$. Let unit $i$ carry $\varepsilon^{\rm o}_i$ and $\varepsilon^{\rm h}_i$ expected errors of each kind, with means $\varepsilon^{\rm o}$ and $\varepsilon^{\rm h}$, and let $\mathcal R$ be the reviewed units. Expected undetected errors are $\sum_{i\in\mathcal R}[\varepsilon^{\rm o}_i e^{\rm o}+\varepsilon^{\rm h}_ie^{\rm h}]+\sum_{i\notin\mathcal R}(\varepsilon^{\rm o}_i+\varepsilon^{\rm h}_i)\ge\sum_i[\varepsilon^{\rm o}_ie^{\rm o}+\varepsilon^{\rm h}_ie^{\rm h}]$, with $e^{\rm o},e^{\rm h}\le1$ the escape probabilities and equality when every unit is reviewed. Per unit the bound is $\bar u$; multiplying by the loss per error, $L_{\rm err}$ times the unit's cost, and summing over units whose cost is $W_A$ gives the loss bound. Both escape probabilities are nonincreasing in $n_r$ and tend to $\chi$, so $\bar u$ falls to $\chi(\varepsilon^{\rm o}+\varepsilon^{\rm h})$; at $n_r=0$ they equal one. The requirement is met at depth $n_r$ if and only if $\bar u\le\bar\varepsilon$; since $\bar u$ decreases to $\chi(\varepsilon^{\rm o}+\varepsilon^{\rm h})$, some depth meets it if that limit is below $\bar\varepsilon$, and none does if it is above.

(d) For a nonnegative random variable,
\[
\mathbb E\min\{Y,\bar Y\}=\int_0^{\bar Y}P(Y>y)\,dy=1+\int_1^{\bar Y}y^{-p_Y}dy,
\]
which is $1+(1-\bar Y^{1-p_Y})/(p_Y-1)$ for $p_Y\ne1$ and $1+\ln\bar Y$ for $p_Y=1$. As $\bar Y\to\infty$ it tends to $1+1/(p_Y-1)=p_Y/(p_Y-1)$ if $p_Y>1$ and diverges if $p_Y\le1$. By independence the expected loss per unit is $\varepsilon^{\rm m}(1-d^{\rm m})L_{\rm mis}\mathbb E\min\{Y,\bar Y\}$ times the unit's cost; summing over units gives the expression.

(e) Events at provider $m$ form a Poisson process of rate $\nu_{\rm cm}$ (the superposition of outages at rate $\nu_{\rm out}$ and of the regressions that change control does not stop, a thinning of rate $(1-q_{\rm chg})\nu_{\rm reg}$), and each disrupts the loss rate $w_mW_{\rm cm}$, at a cost of
\[
T_{\rm con}\bigl[w_mW_{\rm cm}+w_m^2W_{\rm cm}^2/(2Y_{\rm fb})\bigr].
\]
Summing expected annual costs over providers and using $\sum_mw_m=1$ gives $\nu_{\rm cm}T_{\rm con}[W_{\rm cm}+H_{\rm cm}W_{\rm cm}^2/(2Y_{\rm fb})]$. Since $\sum_mw_m^2\le(\sum_mw_m)^2=1$, with equality if and only if one provider carries the whole exposure, the cost is largest with a single provider; it is nonincreasing in $q_{\rm chg}$ because $\nu_{\rm cm}$ is. \qed

\subsection*{Proof of \autoref{prop:attainable}}

Two facts are used throughout. First, $\mathcal Z_n$ is closed under taking subsets: removing processes from a wave keeps it within capacity because uses are nonnegative. Second, $\mathcal Z_n\subseteq\mathcal Z_{n'}$ for $n\le n'$, because empty waves are allowed.

(a) If $t\le t'$, then $n(t)\le n(t')$, so $\mathcal Z_{n(t)}\subseteq\mathcal Z_{n(t')}$ and $\mathcal A_t\subseteq\mathcal A_{t'}$. Raising a capacity or lowering a use keeps every wave within capacity, so $\mathcal Z_n$ weakly grows; raising $T_{\mathrm{emb}}$ weakly lowers $n(t)$. A union over a larger family of sets is larger. On each $\mathcal F(\vartheta_{\mathcal Z})$ the linear score $R_X$ attains its maximum $R^*_X(\vartheta_{\mathcal Z})$ (\autoref{prop:score}(a) and \autoref{rem:cap}), and the family is finite, so the supremum over $\mathcal A_t$ is $\max_{\mathcal Z\in\mathcal Z_{n(t)}}R^*_X(\vartheta_{\mathcal Z})$, and set inclusion gives its monotonicity. For $t<T_{\mathrm{emb}}$, $n(t)=0$, $\mathcal Z_0=\{\emptyset\}$ and $\mathcal A_t=\mathcal F(\vartheta_0)$. Finally, $\mathcal Z_n$ is contained in the set of all subsets of $\mathcal P$.

(b) By the first fact, $\mathcal Z_n$ contains every subset of $\mathcal P$ if and only if $\mathcal P\in\mathcal Z_n$, that is, if and only if $n\ge\bar n$. Since $n(t)\ge\bar n$ if and only if $t\ge T_{\mathrm{emb}}+(\bar n-1)T_{\mathrm{par}}$, $\mathcal A_t=\mathcal A_\infty$ from that date on. If every process fits a wave on its own, one process per wave gives a split into $N$ waves; if the total uses fit in one wave, $\bar n=1$, and $n(T_{\mathrm{emb}})=1$.

(c) Summing the wave constraints over the at most $n(t)$ waves of a split of $\mathcal Z$ gives $\sum_{p\in\mathcal Z}c^i_p\le n(t)\Gamma_i$. With identical uses, a wave can hold $j$ processes if and only if $jc^i\le\Gamma_i$ for every $i$, that is, $j\le n_{\mathrm{cap}}$, so $|\mathcal Z|\le n(t)n_{\mathrm{cap}}$. Embedding all $N$ processes requires $n(t)\ge\lceil N/n_{\mathrm{cap}}\rceil$, that is, $t\ge T_{\mathrm{emb}}+(\lceil N/n_{\mathrm{cap}}\rceil-1)T_{\mathrm{par}}$. Because $T_{\mathrm{emb}}\ge T_{\mathrm{par}}$ and $n_{\mathrm{cap}}\le\Gamma_i/c^i$, this date is at least $\lceil N/n_{\mathrm{cap}}\rceil T_{\mathrm{par}}\ge NT_{\mathrm{par}}/n_{\mathrm{cap}}\ge Nc^iT_{\mathrm{par}}/\Gamma_i$ for every $i$.

(d) By \autoref{prop:score}(b), (c) and (e), $R^*_X=m_0+(X-k)G+\Pi$, where $\Pi$ depends on $\lambda h$ and $X-k$ but not on $m_0$ or $G$. With $k$ and $\lambda$ fixed, $R^*_X(\vartheta_{\mathcal Z})-R^*_X(\vartheta_0)=\sum_{p\in\mathcal Z}\Delta R_p$. For $\mathcal Z\in\mathcal Z_{n(t)}$, the indicator vector of $\mathcal Z$ is feasible for the linear program that defines $\mathcal V_i(n(t)\Gamma_i)$, by (c), so $R^*_X(\vartheta_{\mathcal Z})-R^*_X(\vartheta_0)\le\mathcal V_i(n(t)\Gamma_i)$; maximizing over $\mathcal Z$ gives the bound. The value of a linear maximization as a function of its right-hand side is concave and piecewise linear, and it is nondecreasing here because the feasible set grows with $C$. An optimal solution uses no process with $\Delta R_p\le0$, takes every process with $c^i_p=0$ and $\Delta R_p>0$, and otherwise fills capacity in decreasing order of $\Delta R_p/c^i_p$: if $x_j>0$ and $x_l<1$ with $\Delta R_l/c^i_l>\Delta R_j/c^i_j$, moving $\varepsilon c^i_l$ units of capacity from $j$ to $l$ raises the objective by $\varepsilon c^i_l(\Delta R_l/c^i_l-\Delta R_j/c^i_j)>0$. \qed

\subsection*{Proof of \autoref{prop:expected}}

Let $\mathcal B:=\{\vartheta_0+\sum_px_p\delta\vartheta_p:x\in[0,1]^N\}$. It is compact, $\vartheta_t\in\mathcal B$ always, and $R^*_X$ is continuously differentiable on $\mathcal B$ (\autoref{prop:envelope}), hence bounded, so every expectation below exists. Plans may randomize. Write $\mathcal Z^{\rm e}_t$ for the set of processes embedded by $t$. Three observations are used. (i) Reverting a process removes its increment, and the firm may revert at any time. (ii) Under the uniform clock, the model with freezes at calendar date $t$ is the model without freezes at program time $t_R$ in which disruptions arrive at the rate $\nu_D(1+\nu^{\rm R}T_{\rm frz})$ per unit of program time, because a unit of program time lasts $1+\nu^{\rm R}T_{\rm frz}$ units of calendar time; in particular an attempt survives its flight with probability $\exp\{-\nu_D(1+\nu^{\rm R}T_{\rm frz})T_{\mathrm{emb}}\}$. (iii) Drawing every share $b_p$ at date 0 and revealing it to the plan only when $p$ first embeds does not change the law of anything the plan observes, so we may treat $(b_p)_{p\in\mathcal P}$ as primitive and independent of all other draws.

(a) The plan that never attempts anything gives $\vartheta_t=\vartheta_0$. For the upper bound, fix a plan and a realization. By (ii), each $p\in\mathcal Z^{\rm e}_t$ has a successful attempt that started in a wave $w$ with $(w-1)T_{\mathrm{par}}+T_{\mathrm{emb}}\le t_R$, that is, $w\le n(t_R)$. The successful attempts of each wave are a subset of that wave's attempts, which are within capacity, so $\mathcal Z^{\rm e}_t\in\mathcal Z_{n(t_R)}$. Now $\vartheta_t=\vartheta_0+\sum_{p\in\mathcal Z^{\rm e}_t}b_p\delta\vartheta_p$. Raising the shares $b_p$ to one, one at a time, moves $\vartheta$ along the directions $\delta\vartheta_p$ inside $\mathcal B$, and along each such segment $R^*_X$ has derivative $\nabla R^*_X\cdot\delta\vartheta_p\ge0$ for $X$-improving processes. Hence $R^*_X(\vartheta_t)\le R^*_X(\vartheta_{\mathcal Z^{\rm e}_t})\le R^{\mathrm{att}}_X(t_R)$. Taking expectations and the supremum over plans gives the bound.

(b) \emph{Time.} Given a plan for the horizon $t$ and a later date $t'$, follow the plan until $t$, make no attempt after $t$, and revert each process at the moment it would embed after $t$, by (i). Then $\vartheta_{t'}=\vartheta_t$, so $\bar R_X(t')\ge\mathbb E\,R^*_X(\vartheta_t)$ for every plan, and $\bar R_X(t')\ge\bar R_X(t)$.

\emph{Capacities and uses.} The set of plans grows when a capacity rises or a use falls.

\emph{Success.} Let $\varpi'_p>\varpi_p$, and take any plan of the model with $\varpi_p$. In the model with $\varpi'_p$, run the same plan, except that whenever an attempt of $p$ works, an independent coin that comes up heads with probability $\varpi_p/\varpi'_p$ is drawn; on tails, $p$ is reverted when the outcome is revealed, and the plan is told that the attempt failed. Each attempt of $p$ then works with probability $\varpi_p$, independently of everything else, so the joint law of the plan's information and of $\vartheta_t$ is the same as in the model with $\varpi_p$. Hence the value with $\varpi'_p$ is at least the value with $\varpi_p$.

\emph{Disruption.} Let $\nu'_D>\nu_D$, and take any plan of the model with $\nu'_D$. In the model with $\nu_D$, superpose on the calendar flight of every attempt an independent Poisson process of rate $\nu'_D-\nu_D$, revert the process at its first point, and report that point to the plan as a disruption. The first point of the superposed processes arrives at rate $\nu'_D$, so the joint law is that of the model with $\nu'_D$, and the value with $\nu_D$ is at least the value with $\nu'_D$. The parameters $\nu^{\rm L}$, $\nu^{\rm V}$ and $p_{\rm keep}$ enter only through $\nu_D$ (and, for a sponsorship reset, through $\varpi_p$, which the previous step covers).

\emph{Freezes.} By (ii), $\bar R_X(t)$ equals the value at program time $t_R$ of the model without freezes with disruption rate $\nu_D(1+\nu^{\rm R}T_{\rm frz})$. That value is nondecreasing in its time argument (the first step) and nonincreasing in its disruption rate (the previous step), and raising $\nu^{\rm R}T_{\rm frz}$ lowers $t_R$ and raises the disruption rate.

\emph{Embedding time.} Let $T'_{\mathrm{emb}}>T_{\mathrm{emb}}$ and let every process be $X$-improving. Take any plan of the model with $T'_{\mathrm{emb}}$. In the model with $T_{\mathrm{emb}}$, run it with the same waves, and add to the real disruption process, which covers the calendar flight of an attempt in the model with $T_{\mathrm{emb}}$, an independent Poisson process of rate $\nu_D$ on the further calendar time $(1+\nu^{\rm R}T_{\rm frz})(T'_{\mathrm{emb}}-T_{\mathrm{emb}})$ that the longer flight would last, reverting the process at its first point. The combined disruption process is Poisson with rate $\nu_D$ over the whole longer flight, so the plan's information has the law of the model with $T'_{\mathrm{emb}}$. At every date, every process that is embedded in that model (started at program time $t_0$ with $t_0+T'_{\mathrm{emb}}\le t_R$, working and not disrupted) is also embedded in the model with $T_{\mathrm{emb}}$, with the same share; the additional processes embedded only in the faster model have increments $b_p\delta\vartheta_p$ with $b_p\in[0,1]$. Adding them moves $\vartheta_t$ along improving directions inside $\mathcal B$, so $R^*_X$ does not fall, and the value with $T_{\mathrm{emb}}$ is at least the value with $T'_{\mathrm{emb}}$.

\emph{Realized shares.} Let every process be $X$-improving, and let $b'_p$ have distribution $F'_p\le F_p$, where $F_p$ is the distribution of $b_p$. Take any plan of the model with $F_p$. When $p$ first embeds in the model with $F'_p$, draw an independent uniform $U$, let $\tilde U:=F'_p(b'_p-)+U[F'_p(b'_p)-F'_p(b'_p-)]$, which is uniform on $[0,1]$ with $F'^{-1}_p(\tilde U)=b'_p$ almost surely, and report $\tilde b_p:=F^{-1}_p(\tilde U)$ to the plan. Then $\tilde b_p$ has distribution $F_p$, and $\tilde b_p\le b'_p$ because $F_p\ge F'_p$ implies $F^{-1}_p\le F'^{-1}_p$. The plan's actions have the law of the model with $F_p$, and the realized parameters exceed that model's in every share; later embeddings of $p$ reuse the same shares in both models. Since $R^*_X$ is nondecreasing in each share on $\mathcal B$, the value with $F'_p$ is at least the value with $F_p$.

{\sloppy (c) Let $\mathcal S$ be the set of attempts that start in waves $1,\dots,n(t_R)$. As in (a), every process embedded by $t$ has a successful, undisrupted attempt in $\mathcal S$, and distinct processes have distinct attempts, so $|\mathcal Z^{\rm e}_t|\le\sum_{\iota\in\mathcal S}\mathbf 1\{\iota\text{ works and is not disrupted}\}$. The decision to make an attempt depends only on information available when it starts, and its outcome is independent of that information and works and survives with probability $p^{\rm emb}$ by (ii). Summing over the at most $n_{\mathrm{cap}}$ attempts of each of the $n(t_R)$ waves, $\mathbb E|\mathcal Z^{\rm e}_t|\le p^{\rm emb}\,\mathbb E|\mathcal S|\le p^{\rm emb}\,n_{\mathrm{cap}}n(t_R)$, and $|\mathcal Z^{\rm e}_t|\le N$. For equality, let the plan attempt, in each wave, $n_{\mathrm{cap}}$ processes that are neither embedded nor in flight, and never revert a process voluntarily. At the start of wave $w\le n(t_R)$ at most $n_{\mathrm{cap}}(w-1)\le N-n_{\mathrm{cap}}$ processes are embedded or in flight when $N\ge n_{\mathrm{cap}}n(t_R)$, so the wave can be filled; a process is attempted again only after a failure or a disruption, so successes belong to distinct processes, and $\mathbb E|\mathcal Z^{\rm e}_t|=p^{\rm emb}\,n_{\mathrm{cap}}n(t_R)$. For the score, \autoref{prop:attainable}(d) gives $R^*_X(\vartheta_t)=R^*_X(\vartheta_0)+\Delta R\sum_{p\in\mathcal Z^{\rm e}_t}b_p$. Let $\mathcal F_t$ be the set of processes that have embedded at least once by $t$; then $\mathcal Z^{\rm e}_t\subseteq\mathcal F_t$, and $|\mathcal F_t|$ is at most the number of successful undisrupted attempts in $\mathcal S$. By (iii), $b_p$ is independent of whether and when $p$ first embeds, so $\mathbb E\sum_{p\in\mathcal Z^{\rm e}_t}b_p\le\mathbb E\sum_{p\in\mathcal F_t}b_p=\bar b\,\mathbb E|\mathcal F_t|\le\bar b\,p^{\rm emb}\,n_{\mathrm{cap}}n(t_R)$, and $\sum_{p\in\mathcal Z^{\rm e}_t}b_p\le N$. Because a process realizes the same share at every embedding, reverting a process and migrating it again cannot raise its share, so the bound holds for every plan.\par}

(d) Let the plan attempt, in each wave and in a fixed cyclic order, as many processes that are neither embedded nor in flight as capacity allows, which is at least one. Every process that is not embedded is then attempted infinitely often until it embeds, and each attempt embeds it with probability $p^{\rm emb}_p>0$, independently of the past, so the number of attempts it needs is geometric and almost surely finite. After the last process embeds, $\vartheta_t=\vartheta_\infty:=\vartheta_0+\sum_pb_p\delta\vartheta_p$. Hence $R^*_X(\vartheta_t)\to R^*_X(\vartheta_\infty)$ almost surely, and bounded convergence gives the limit of expectations. Now let every process be $X$-improving, and fix any plan and date. With the shares of (iii), $\vartheta_t=\vartheta_0+\sum_{p\in\mathcal Z^{\rm e}_t}b_p\delta\vartheta_p$, and adding $b_p\delta\vartheta_p$ for every $p\notin\mathcal Z^{\rm e}_t$ moves $\vartheta$ along improving directions inside $\mathcal B$, so $R^*_X(\vartheta_t)\le R^*_X(\vartheta_\infty)\le R^*_X(\vartheta_{\mathcal P})$, the second inequality by raising every share to one as in (a). Taking expectations and the supremum over plans, $\bar R_X(t)\le\mathbb E\,R^*_X(\vartheta_\infty)$, and the retry plan attains this in the limit. In the setting of (c), $\mathbb E\,R^*_X(\vartheta_\infty)=R^*_X(\vartheta_0)+N\bar b\,\Delta R=R^*_X(\vartheta_{\mathcal P})-(1-\bar b)N\Delta R$.

For the last claim, let $\mathcal Z^\star$ maximize $R^*_X(\vartheta_{\mathcal Z})$ over all subsets, so that $R^*_X(\vartheta_{\mathcal Z^\star})=R^{\mathrm{att}}_X(\infty)$. Once $\Gamma_i\ge\sum_{p\in\mathcal Z^\star}c^i_p$ for every $i$, the plan that attempts $\mathcal Z^\star$ in the first wave is feasible. Fix $t>T_{\mathrm{emb}}$; once $\nu^{\rm R}T_{\rm frz}$ is small enough, $t_R\ge T_{\mathrm{emb}}$, and the plan embeds all of $\mathcal Z^\star$ by $t$ with probability $\prod_{p\in\mathcal Z^\star}p^{\rm emb}_p\to1$. On that event $\vartheta_t=\vartheta_0+\sum_{p\in\mathcal Z^\star}b_p\delta\vartheta_p$, which converges in probability to $\vartheta_{\mathcal Z^\star}$. Because $R^*_X$ is continuous and bounded on $\mathcal B$, $\mathbb E\,R^*_X(\vartheta_t)\to R^{\mathrm{att}}_X(\infty)$, so $\liminf\bar R_X(t)\ge R^{\mathrm{att}}_X(\infty)$. By (a) and \autoref{prop:attainable}(a), $\bar R_X(t)\le R^{\mathrm{att}}_X(t_R)\le R^{\mathrm{att}}_X(\infty)$, because every process is $X$-improving.

(e) Measure the program in units of mass, and let $A_w$ be the mass attempted in waves $1,\dots,w$. Each wave attempts at most $c_w$, so $A_n\le c_wn$. Of the mass attempted in waves $1,\dots,n(t_R)$, the fraction $p^{\rm emb}$ embeds by $t$, and nothing attempted later does, so $\mathcal E(t)\le p^{\rm emb}A_{n(t_R)}\le p^{\rm emb}c_wn(t_R)$, and $\mathcal E(t)\le1$. By \autoref{prop:attainable}(d) applied to the mass, $R^*_X(\vartheta_t)=R^*_X(\vartheta_0)+\Delta\hat R$ times the realized benefit of the embedded mass, and by (iii) the shares average to $\bar b$ over any embedded mass, which gives (i). For (ii), at the start of wave $w$ the mass that is embedded or working in flight is at most $\varpi A_{w-1}$, because failed mass is revealed and reverts by the end of its wave; so at least $1-\varpi A_{w-1}$ is available. The greedy plan has $A_{w-1}=c_w(w-1)$ as long as it has filled every earlier wave, and it can fill wave $w$ if $1-\varpi c_w(w-1)\ge c_w$. If $c_w[1+\varpi(n-1)]\le1$, this holds for every $w\le n$, so waves $1,\dots,n$ are filled and $\mathcal E(t)=p^{\rm emb}c_wn$ for $n=n(t_R)$. For (iii), with $\nu_D=0$ working mass is never disrupted, so the available mass at the start of wave $w$ is exactly $U_{w-1}$, where $U_w:=1-\varpi A_w$ and $U_0=1$, and $U_w=U_{w-1}-\varpi\min\{c_w,U_{w-1}\}$; all working mass from waves $1,\dots,n(t_R)$ has embedded by $t$, so $\mathcal E(t)=1-U_{n(t_R)}$. While $U_{w-1}\ge c_w$, $U_w=1-\varpi c_ww$ and $\mathcal E=\varpi c_wn$. Let $w_1$ be the first wave with $U_{w_1-1}<c_w$; from then on $U_w=(1-\varpi)U_{w-1}<c_w$, so $\mathcal E(t)>1-c_w\ge\min\{1,\varpi c_wn(t_R)\}-c_w$, and $U_w\to0$. Since $\mathcal E(t)\le\min\{1,\varpi c_wn(t_R)\}$ by (i), the shortfall is less than $c_w$.

(f) Let $N^{\rm R}(u)$ be the number of reorganizations in program time $[0,u)$, a Poisson process of rate $\nu^{\rm R}$. The calendar time at which the program reaches program time $u$ is $C(u)=u+T_{\rm frz}N^{\rm R}(u)$, and $\tilde t(t)=\sup\{u:C(u)\le t\}$. By the strong law of large numbers $N^{\rm R}(u)/u\to\nu^{\rm R}$ almost surely, so $C(u)/u\to1+\nu^{\rm R}T_{\rm frz}$; since $C(\tilde t(t)-)\le t\le C(\tilde t(t))$ and $\tilde t(t)\to\infty$, $\tilde t(t)/t\to1/(1+\nu^{\rm R}T_{\rm frz})$. For the score bound, the argument of (a) applies path by path with $\tilde t(t)$ in place of $t_R$: every process embedded by calendar date $t$ has a successful attempt in a wave $w\le n(\tilde t(t))$, so $R^*_X(\vartheta_t)\le R^{\mathrm{att}}_X(\tilde t(t))$, and taking expectations gives the bound. For the count, condition on the whole freeze path, which is independent of everything else. Given it, an attempt started at program time $t_0$ has a calendar flight of length $L=T_{\mathrm{emb}}+T_{\rm frz}[N^{\rm R}(t_0+T_{\mathrm{emb}})-N^{\rm R}(t_0)]\ge T_{\mathrm{emb}}$, and it works and survives with probability $\varpi e^{-\nu_DL}\le\varpi e^{-\nu_DT_{\mathrm{emb}}}$, independently of the information at its start; the set of attempts that start in waves $1,\dots,n(\tilde t(t))$ is known at their start given the freeze path. The argument of (c) then gives the conditional bound $\varpi e^{-\nu_DT_{\mathrm{emb}}}n_{\mathrm{cap}}n(\tilde t(t))$, and taking expectations gives the stated one. Unconditionally, the number of reorganizations during a flight is Poisson with mean $\nu^{\rm R}T_{\mathrm{emb}}$, so the survival probability is $e^{-\nu_DT_{\mathrm{emb}}}\,\mathbb E\,e^{-\nu_DT_{\rm frz}N}=\exp\{-\nu_DT_{\mathrm{emb}}-\nu^{\rm R}T_{\mathrm{emb}}(1-e^{-\nu_DT_{\rm frz}})\}$, at least the uniform-clock value because $1-e^{-x}\le x$. Finally, take one wave that holds the whole program, $T_{\mathrm{emb}}=T_{\mathrm{par}}=1$, $\varpi=1$, $\nu_D=0$, $\nu^{\rm R}=0.5$ and $T_{\rm frz}=2$. At $t=1.5$ the uniform clock gives $t_R=0.75<T_{\mathrm{emb}}$, so $R^{\mathrm{att}}_X(t_R)=R^*_X(\vartheta_0)$ and the bound of (c) is zero. With random freezes, no reorganization occurs in the first unit of program time with probability $e^{-0.5}\approx0.61$, and then the program embeds at calendar time 1; so the expected number embedded at $t=1.5$ is at least $0.61N$, and $\bar R_X(1.5)$ exceeds $R^*_X(\vartheta_0)$ whenever the program raises the score. \qed

\subsection*{Proof of \autoref{lem:expected-frontier}}

(a) For each realization, $m^*(\cdot;\vartheta_t)$ is concave on $[G(\vartheta_t),\infty)$ (\autoref{prop:frontier}(c)) and equal to $-\infty$ to the left of it, so it is concave on the real line as an extended-valued function. It is bounded above on $\mathcal B$, so its expectation is well defined and, as an average of concave functions, concave. If only $m_0$ and $k$ vary, $G$ is fixed and $m^*(g;\vartheta)=m_0-kg-\kappa(g-G)^{\beta}$ is affine in $(m_0,k)$ for each $g$, so $\mathbb E\,m^*(g;\vartheta_t)=m^*(g;\mathbb E\vartheta_t)$. (b) For every $g$ and realization, $Xg+m^*(g;\vartheta_t)\le R^*_X(\vartheta_t)$ by \autoref{prop:score}(a). Taking expectations gives $Xg+m^{\mathbb E}_t(g)\le\mathbb E\,R^*_X(\vartheta_t)$ for every $g$, and a firm that fixes $g$ in advance and then earns the frontier margin at $g$ scores $Xg+m^{\mathbb E}_t(g)$ in expectation. \qed

\subsection*{Proof of \autoref{rem:regimes}}

(i) The attainable set is a finite union of sets on each of which $R_X$ attains its maximum, so maximizers exist. Let $(g,m)$ maximize $R_X$ and $(g',m')$ maximize $R_{X'}$ on it, with $X<X'$. Then $Xg+m\ge Xg'+m'$ and $X'g'+m'\ge X'g+m$; adding gives $(X'-X)(g'-g)\ge0$, so $g'\ge g$. If $g'>g$, $m'\ge m$ and $X>0$, then $Xg'+m'>Xg+m$, which contradicts the optimality of $(g,m)$; so $m'<m$. (ii) Every split into waves within capacity of a subset of $\mathcal P'$ is also one of a subset of $\mathcal P$, so the families, and with them the attainable and static sets, are nested. (iii) A tighter dip budget lowers the dip capacity, and by hypothesis slower decisions lower the decision capacity; Propositions~\ref{prop:attainable}(a) and \ref{prop:expected}(b) give the effect on $\mathcal A_t$ and $\bar R_X(t)$. The static set $\mathcal A_\infty$ does not depend on capacities. \qed

\subsection*{Proof of \autoref{prop:general}}

(a) Substituting the frontier \eqref{eq:frontier} into $e=m+k(g-r)$ gives $e^*(g)=m_0-rk-\kappa(g-G)^{\beta}$. Its derivative $-\beta\kappa(g-G)^{\beta-1}$ tends to zero as $g\downarrow G$ and is negative for $g>G$, and its second derivative is that of $m^*$, so it is concave by \autoref{prop:frontier}(c). For $k>0$, $(m^*+kg)/k=[m_0-\kappa(g-G)^{\beta}]/k$. For each $g$ the maps $m\mapsto m+k(g-r)$ and $m\mapsto(m+kg)/k$ are strictly increasing, so the largest feasible $m$ at $g$, produced by the spend $\kappa(g-G)^{\beta}$ (\autoref{prop:frontier}(a)), gives the largest feasible economic margin and ROIC at $g$. Both frontiers are strictly decreasing for $g>G$, so every point on them is undominated.

(b) $(X-k)g+m+k(g-r)+rk=Xg+m$, and $k[(m+kg)/k+(X/k-1)g]=m+kg+Xg-kg=Xg+m$. Since $q_K>1$ if and only if $X>k$, the last claim is \autoref{prop:score}(b) and (c).

(c) With revenue $Se^{gt}$ and a constant margin, $V=\int_0^\infty e^{-rt}mSe^{gt}\,dt=mS/(r-g)$ for $g<r$. Substituting $m=e-k(g-r)$ gives $V/S=k+e/(r-g)$, and dividing by $k$ and substituting $e=k(\mathrm{ROIC}-r)$ gives $V/(kS)=1+(\mathrm{ROIC}-r)/(r-g)=(\mathrm{ROIC}-g)/(r-g)$. The FCF yield is $mS/V=r-g$. The set $V/S=v$ is the line $e=(v-k)(r-g)$, which passes through $(r,0)$ with slope $-(v-k)$. By \autoref{prop:dynamic}(d), $v_0=m^\circ/(r-g^\circ)$, and $(g^\circ,m^\circ)$ maximizes $R_{v_0}$ on the frontier with $\ell=\lambda$; if $v_0>k$ the optimum is interior and $-m^{*\prime}(g^\circ)=v_0$ (\autoref{prop:score}(b)). Hence $-e^{*\prime}(g^\circ)=-m^{*\prime}(g^\circ)-k=v_0-k$ and $e^*(g^\circ)=m^\circ+k(g^\circ-r)=(v_0-k)(r-g^\circ)$.

(d) With $k=0$, $e=m$ by definition. With $X=1$, (b) gives $R_1=(1-k)g+e+rk$, whose derivative in $g$ at fixed $(e,k)$ is $1-k$ and in $k$ at fixed $(g,e)$ is $r-g$. In the example, $(r-g)\times1.5=-0.075$ and $(1-1.5)\times0.01=-0.005$. \qed

\subsection*{Proof of \autoref{prop:levers-position}}
By \eqref{eq:sigstar} and \eqref{eq:Rstar}, $\sigma^*=(\alpha\ell y)^{1/(1-\alpha)}$, $y\Delta^*=\sigma^*/\alpha$ and $\Pi=C_\alpha(\ell y)^{1/(1-\alpha)}$, with $y=X-k$. By \autoref{prop:envelope}, a change $d\ln\lambda$ raises $R^*_X$ by $y\Delta^*d\ln\lambda=(\sigma^*/\alpha)d\ln\lambda$, which is $x\sigma^*/\alpha$ for a gain $x$ in demand pull and, to first order, $x\sigma^*$ for the saving $x$, since $d\ln\lambda=-\alpha\ln(1-x)$. A cut $xa$ raises $m_0$, and $x(1-\epsilon_p)c_h$ raises $\gamma$ (\autoref{lem:prices}), each one for one in $R^*_X$ and in $R_\infty=m_0+(X-k)G_i$; $\lambda$ does not enter $R_\infty$ and $\Pi\to0$ as $s\to\bar s$ (\autoref{prop:saturation}). Retention and capital are \autoref{cor:priorities}(iii) and \autoref{cor:capital} with \autoref{rem:assumptions}(b), (c). This proves (a). For (b), $x\sigma^*>xa$ if and only if $\alpha\ell y>a^{1-\alpha}$, and $x\sigma^*/\alpha>xa$ if and only if $\alpha\ell y>(\alpha a)^{1-\alpha}$; if $\hat\sigma=(\alpha\ell(\hat X-k))^{1/(1-\alpha)}$ then $\alpha\ell=\hat\sigma^{1-\alpha}/(\hat X-k)$. For (c), $\ell=\lambda h(s)$ and $\sigma^*$ is proportional to $\ell^{1/(1-\alpha)}$. Part (d) is \autoref{cor:priorities}(iv). For (e), $m_0$ and $G_i$ enter \eqref{eq:Rstar} only through $m_0+(X-k)G_i$. For (f), $\sigma^*$ increases in $X$ and $xa$ and $x(1-\epsilon_p)c_h$ do not depend on it.

\subsection*{Proof of \autoref{rem:owncompute}}

Bought as a service, the compute adds $(1+\upsilon_c)\mathrm{TCO}$ to $c_c$, which costs $(1-\tau)(1+\upsilon_c)\mathrm{TCO}$ of the after-tax score. Owned, it adds $(1-s_{\rm srv})\mathrm{TCO}$ to $c_c$, the replacement spending $k_c/L_{\rm srv}=s_{\rm srv}\mathrm{TCO}$ and the capital $k_c=s_{\rm srv}L_{\rm srv}\mathrm{TCO}$. By \autoref{cor:ceiling} with $X=1$, the running cost costs $(1-\tau)(1-s_{\rm srv})\mathrm{TCO}$, replacement $(1-\psi\tau)s_{\rm srv}\mathrm{TCO}$ and growth capital $(1-\psi\tau)k_cG$. Owning minus buying therefore changes the ceiling by $(1-\tau)\upsilon_c\mathrm{TCO}+[(1-\tau)-(1-\psi\tau)]s_{\rm srv}\mathrm{TCO}-(1-\psi\tau)s_{\rm srv}L_{\rm srv}\mathrm{TCO}\,G$, and $(1-\tau)-(1-\psi\tau)=-(1-\psi)\tau$. At $\psi=1$ the change is $(1-\tau)[\upsilon_c-s_{\rm srv}L_{\rm srv}G]\mathrm{TCO}$, positive if and only if $\upsilon_c>s_{\rm srv}L_{\rm srv}G$; with $s_{\rm srv}L_{\rm srv}=0.6\times5.5=3.3$ the threshold is $0.165$ at $G=0.05$ and $0.495$ at $G=0.15$. At $\psi=0$ the threshold is $[\tau s_{\rm srv}+s_{\rm srv}L_{\rm srv}G]/(1-\tau)$, which at $\tau=0.18$ is $0.333$ and $0.735$. \qed

\subsection*{Proof of \autoref{rem:assumptions}}

(a) At a given date $S$ is given, so $m_0=\gamma-a_v-\Xi/S$ is a constant in the static problem and Lemmas~\ref{prop:frontier}, \ref{prop:score} and \ref{prop:envelope} apply as stated. $R_\infty=\gamma-a_v-\Xi/(sM)+(X-k)G$ is increasing in $S=sM$. At fixed $M$, $R_X^*(s)=R_\infty(s)+\Pi(s)$, and by \autoref{prop:saturation}(a) $d\Pi/ds=-(X-k)_+\Delta^*/(\bar s-s)$, which is bounded as $s\to 0$, while $d R_\infty/ds=\Xi/(s^2M)\to\infty$; so $R_X^*$ is increasing at small $s$. Along a path with largest revenue $S_{\max}$, $m_t\le m_0(S_t)-\sigma_t-kg_t\le m_0(S_{\max})-\sigma_t-kg_t$, so the proof of \autoref{prop:budget} applies with $m_0(S_{\max})$. For \autoref{prop:dynamic}, the FCF flow is $(\gamma-a_v-\sigma-kg)S-\Xi$, so for $r>0$ the value of any policy is its value with intercept $\gamma-a_v$ minus $\int_0^\infty e^{-rt}\Xi\,dt=\Xi/r$, which does not depend on the policy; the optimal policy and $\partial V/\partial S$ are therefore those of \autoref{prop:dynamic} with $m_0=\gamma-a_v$.

(b) With $g=G_i+\ell\sigma^{\alpha}$, the static problem at a date is that of Section~\ref{sec:model} with $G$ replaced by $G_i$, so the static propositions hold with this substitution. Along any path, $R_X\le Xg+m_0-kg=m_0+(X-k)G_i+(X-k)(g-G_i)$, and $g\ge G_i$, so $(R_X-R_\infty)^+\le(X-k)(g-G_i)$ with $R_\infty=m_0+(X-k)G_i$. Since $\dot s/s=g-G$, $\int_0^T(g-G_i)\,dt=\ln(s(T)/s_0)+(G-G_i)T$.

(c) For $k<0$, the proofs of \autoref{prop:score}(a),(b),(e), \autoref{prop:saturation}, \autoref{prop:envelope} and \autoref{prop:budget} do not use the sign of $k$ (for $X\ge 0>k$, $y>0$, so part (b) of \autoref{prop:score} applies). The frontier satisfies $m^{*\prime}(g)=-k-\beta\kappa(g-G)^{\beta-1}$, which is positive for $g-G<\Delta_0$ and negative beyond, so $m^*$ is strictly increasing on $[G,G+\Delta_0]$ and every point with $g<G+\Delta_0$ is dominated by $(G+\Delta_0,m^*(G+\Delta_0))$; strict concavity is unaffected. A single firm with $k<0$ already makes $\Phi=m^*$ increasing on $[G,G+\Delta_0]$, so \autoref{prop:industry}(i) can fail. The proof of \autoref{prop:dynamic}(d) uses $k\ge 0$ in $\mathcal H(0)<0$ and in $m_t\ge m^\circ$. With depreciation, investment is $\dot K+\delta_KK=k(g+\delta_K)S$, so $m=(m_0-\delta_Kk)-\sigma-kg$, which is \eqref{eq:fcf} with intercept $m_0-\delta_Kk$. Then $R_X^*(k)=\mathcal R(k)-\delta_Kk$ with $\mathcal R$ convex as in the proof of \autoref{cor:capital}, so $\partial R_X^*/\partial k=-(g^*+\delta_K)$ and the bracket \eqref{eq:kl-bracket} holds with $g^*$ replaced by $g^*+\delta_K$. At $g^*=0.10$, $\delta_K=0.25$ and $dk=0.10$ the hurdle is $0.035$; at $g^*=0.50$ it is $0.075$.

(d) The proofs of Lemmas~\ref{prop:frontier}, \ref{prop:score} and \ref{prop:envelope} and of Propositions~\ref{prop:industry} and \ref{prop:industry-full} take $\ell$ as given at a point and never use $h'$. The proof of \autoref{prop:budget} uses the functional form only to show $s(t)<\bar s$, which the cap $s\le\bar s$ provides directly. Monotonicity of $h$ enters the proof of \autoref{prop:saturation}(a) (sign of $dR_X^*/ds$) and (c) (the supremum is the limit at $s\to 0$), the attainability statement in (d) (the range of $\lambda h$ is $(0,\lambda)$), the bound $\ell(s)\le\lambda$ and the logistic path in the proofs of \autoref{prop:dynamic}(d) and (e), and the logistic solution in the proof of \autoref{prop:myopic}. \qed
